% 本文件由 LaTeX 排版 Agent 根据有效 translation.json 自动生成。
% author=Tertullian; work_id=0407; title=论端庄

\chapter{论端庄}

% structure: p0001 source=h1 target=omitted reason=page-title-already-rendered-by-parent

% structure: p0002 source=h2 target=section reason=relative-heading-rank; base=h2

\section{第一章. 天主既是公义也是仁慈的；因此，慈悲不可滥施。}

端庄，是举止的花朵，我们身体的荣誉，性别的优雅，血统的完整，民族的保障，圣德的基础，一切良好秉性的预兆；虽然它稀罕，不易臻于完美，而且几乎不能永远保持，但若本性已为其打下初步基础，纪律劝勉它，监察的严厉约束其过度，它在世上仍会在某种程度上存留——即基于这样的假设：每一种心灵善质要么源于出生，要么源于教养，要么源于外在强制。\par

但是，由于邪恶之事的征服力量正在增长——这是末世的特征——善事如今既不允许生发，因为生殖原理已被败坏；也不允许受训，因为学问已荒废；也不得强制，因为法律已被解除武装。事实上，我们现在开始讨论的（端庄）此刻已变得如此过时，以至人们认为端庄不是对欲望的弃绝，而是对欲望的节制；并且，他若没有\Emph{过于}贞洁，就被视为\Emph{足够}贞洁了。但是，让世界的端庄与它自身以及世界本身去处理吧：连同它那若习惯上源于出生的本性，它的若源于教养的学问，它的若源于强制的奴役；只是，如果它仅仅留存下来却证明无效，它本会更不幸，因为它的活动不曾在天主的家庭中施行。我宁愿没有善，也不要有虚浮的善：那存在而无益的存在有什么益处呢？是我们的善事现状正在沉沦；是基督徒端庄的体系正被动摇根基——（基督徒端庄），它的一切都源于天国：它的本性，\Quote{借着重生的圣洗}；它的纪律，借着讲道的媒介；它的监察严厉，借着两约各自展现的审判；并服从于更恒常的外在强制，源于对永火或天国的恐惧或渴望。\par

反对这种（端庄），我岂能装假不言呢？我听说甚至有一道敕令已经颁布，而且是强制性的。\OriginalTerm{最高司祭}{Pontifex Maximus}——即主教中的主教——发布敕令：\Quote{我赦免那些完成了（悔改要求的）人，他们的奸淫罪和邪淫罪。}啊，敕令，上面不能刻上\InnerQuote{善行！}而这宽大又将在何处张贴呢？我想，就在那感官欲望的门口，就在感官欲望的名号之下。那里是宣扬这种悔改的地方，正是过犯自身萦绕之处。那里是宣读赦免的地方，人们将怀此希望进入。但这份（敕令）却在教会中宣读，在教会中宣告；而（教会）是一位贞女！愿这样的宣告远离基督的净配！她，真正的、端庄的、圣洁的，连耳朵也毫无玷污。她没有任何人可许下这样的承诺；即便她曾有，她也不会许下；因为就连天主在地上的圣殿，主也早就可以称之为\BibleQuote{贼窝}，而非奸淫者和邪淫者的窝。\par

因此，这也将是我控告\OriginalTerm{精神派}{Psychics}的罪状之一；也是我先前与他们同感一心的罪状；好让他们更以此来指责我的反复无常。与同道绝交从来不是犯罪的先兆。难道与多数人一同犯错不更容易吗？真理却是与少数人同在！然而，有益的反覆无常不再是我的耻辱，正如我不希望有害的反覆无常成为我的装饰。我不因自己所放弃的错误而羞惭，因为我为自己不再持有它而欣喜，因为我意识到自己变得更好、更稳重。没有人因自己的进步而羞愧。即使在基督内，知识也有成长的阶段；宗徒也经历了这些阶段。\BibleQuote{当我还是孩子的时候，说话像孩子，理解像孩子；但当我成为成年人时，就把孩子的事丢弃了。}他确实如此地放弃了他早先的观点：他成为传统（不是祖传的，而是基督信仰的传统）的效仿者，甚至渴望那些主张保留割礼者的精确性，这并非犯罪。但愿同样的命运也降临到那些割裂肉体纯正真实完整的人身上；他们不是割掉最表面的一层，而是割掉了谦逊本身最深处的形象，同时向奸淫者和行淫者承诺赦免，公然违背基督教之名的首要纪律；这纪律连外邦人本身也如此强调地作证，以至于他们通过肉体的玷污而非折磨来惩罚我们女性中的这纪律，想要从她们那里夺走她们比生命更珍视的东西！但现在这荣耀正在被熄灭，而且是被那些本应更加坚定地拒绝对此类玷污给予任何赦免的人所熄灭，因为他们把害怕屈服于奸淫和行淫作为随意结婚的理由——因为\BibleQuote{与其欲火中烧，不如结婚}。无疑，为了贞洁，不贞是必要的——\Quote{欲火}将被\Quote{火}扑灭！那么，他们为什么还要以悔改为名，对那些他们通过多婚法提供补救的罪行给予宽恕呢？因为如果罪行被宽恕，补救措施就是徒劳的；如果补救措施徒劳，罪行就会继续存在。因此，无论哪种方式，他们都以忧虑和疏忽为儿戏：对他们允许通融的罪行采取最空洞的预防措施，对他们采取预防措施的罪行给予最荒谬的通融；而实际上，在允许通融的地方就不应采取预防，在采取预防的地方就不应允许通融；因为他们采取预防，仿佛不愿某事发生；但给予宽恕，仿佛愿意它发生。然而，如果他们不愿它发生，就不应该给予宽恕；如果他们愿意给予宽恕，就不应该采取预防。因为，奸淫和行淫不会同时被归为中等罪和最大的罪，以便对它们可以同样地采取两种对策——采取预防的忧虑和给予宽恕的安全。但由于它们位于罪行之巅，就没有余地同时给予它们像中等罪那样的宽恕，以及像最大罪那样的预防。但\Emph{我们}却这样防范最大的（或者，如果你愿意，\Emph{最高的}）罪行，即：相信之后，不允许知道甚至第二次婚姻，尽管它通过婚书和嫁妆与奸淫和行淫的行为有所区别；因此，我们极其严格地将再婚者逐出教会，因为他们以纪律的非常规性给护慰者带来了耻辱。我们为奸淫者和行淫者也设定了同样的界限；判决他们流下没有平安的眼泪，并从教会中除了公布他们的耻辱之外，得不到任何更丰厚的回报。\par

% structure: p0007 source=h2 target=section reason=relative-heading-rank; base=h2

\section{第二章 天主同样公义与仁慈；因此，仁慈不可不分对象。}

\Quote{但是，他们说，天主是\InnerQuote{善良}的，\InnerQuote{至善}的，\InnerQuote{慈悲心肠}的，\InnerQuote{怜悯者}，\InnerQuote{充满慈悲心肠}，祂视这比\InnerQuote{一切祭献}更\InnerQuote{贵重}，\InnerQuote{不认为罪人的死亡比他的悔改更有价值}，\InnerQuote{所有人的救主，尤其是信徒的救主}。所以，\InnerQuote{天主的儿女？}也应当\InnerQuote{慈悲心肠}且\InnerQuote{缔造和平}；\InnerQuote{反过来施与，正如基督也曾施与我们一样}；\InnerQuote{不要审判，免得你们受审判}。因为\InnerQuote{他的主人面前，人或站立或跌倒；你是谁，竟敢审判别人的仆人？}\InnerQuote{宽免，你们也将获得宽免}。}他们用这些如此大的荒唐话谄媚天主、迎合自己，软化而非强化纪律，而我们又能拿出多么有力且相反的论据予以反驳——这些论据警告性地将天主的\Quote{严厉}摆在我们面前，并激发我们自己的恒心？因为，虽然天主本性是善的，但祂也是\Quote{公义的}。因为，出于事情的本性，正如祂知道如何\Quote{医治}，祂也知道如何\Quote{击打}；\Quote{缔造和平}，但也\Quote{制造灾祸}；祂更喜爱悔改，但也命令\Person{耶肋米亚}不要为有罪的百姓祈求免除灾祸——\Quote{如果他们禁食}，祂说，\Quote{我必不听他们的恳求}。又说：\Quote{不要为这百姓向我祈祷，不要为他们祷告祈求，因为我必不听他们在我呼求我的时候、在他们患难之时的祈祷}。此外，上文那位将怜悯置于祭献之上的同一位（天主）说：\Quote{不要为这百姓向我祈祷，不要祈求他们能得怜悯，也不要到我面前为他们求情，因为我必不听他们}——当然是指当他们求怜悯、出于悔改而哭泣禁食、并将自己的克苦献给天主的时候。因为天主是\Quote{忌邪的}，是一位不可被藐视戏弄的神——被那些谄媚祂良善的人戏弄——并且，祂虽然\Quote{忍耐的}，却通过\Person{依撒意亚}警告祂忍耐的终结。\Quote{我沉默已久；难道我要永远沉默并忍耐吗？我像临产妇人一样安静；我要起来，使他们干旱。}因为\Quote{火要在祂面前发出，烧灭祂的仇敌；}不仅击打身体，也击打灵魂，打入地狱。此外，主亲自示范祂如何威胁那些审判的人：\Quote{因为你们用什么审判审判人，你们也要受什么审判。}所以祂并没有禁止审判，而是教导如何审判。因此，宗徒也审判，且是在一件奸淫的事上，说\Quote{这样的人必须交给\Person{撒殚}，败坏他的肉体；}同样责备他们因为\InnerQuote{弟兄们}没有被\InnerQuote{在圣人面前受审判}：因为他接着说，\Quote{审判外人于我何干？}\Quote{但是你们要宽恕，好使你们经由天主获得宽恕。}这样被洁净的罪，是一个人可能对弟兄犯下的，而非对天主犯下的。简而言之，我们在祈祷中宣称，我们将宽免我们的欠债者；但是，以这些圣经的权威为基础，将争论的绳索交替拉扯向不同方向，以至于一段圣经似乎拉紧，另一段放松纪律的缰绳——仿佛悬而未决——后一段因宽大而贬低悔改的补救帮助，前一段因严厉而拒绝它，这样做是不合适的。此外，圣经的权威将限定在其自身范围内，互不矛盾。悔改的补救帮助由其自身条件决定，没有无限制的让步；其原因本身也已在命题中预先区分清楚，没有混乱。我们同意，悔改的原因是罪。我们将这些罪分为两种：一些是可宽恕的，一些是不可宽恕的：据此，没有人会怀疑有些罪应受惩罚，有些应受定罪。每一项罪要么通过宽赦，要么通过刑罚来清除：宽赦是惩罚的结果，刑罚是定罪的结果。关于这一区别，我们不仅已经预先引用了圣经中某些对立的段落，一段保留罪，一段宽恕罪；而且\Person{若望}也将教导我们：\Quote{如果有人看见他的弟兄犯了不至于死的罪，他就应当祈求，天主必将生命赐给他；}因为他没有\Quote{犯至于死的罪，}这就可宽恕。\Quote{有一种至于死的罪；我不是说应当为那罪祈求}——这是不可宽恕的。所以，哪里有\Quote{祈求}的效力，那里就有宽恕的效力；哪里没有\Quote{祈求}的效力，那里同样也没有宽恕。根据罪的这一区别，悔改的条件也得以区分。有一种条件可能得到宽赦——即可宽恕之罪的情形；另一种条件绝不可能得到——即不可宽恕之罪的情形。剩下要特别考察的是，关于通奸和淫乱的位置，它们应当被归入哪一类罪。\par

% structure: p0009 source=h2 target=section reason=relative-heading-rank; base=h2

\section{一项在上文承诺的讨论开始前预料到的反对意见}

但在做此事之前，我将简短回答从对立面遇到的一种异议，即关于我们刚刚定义为无宽恕的那种悔改。\Quote{为什么，如果，}他们说道，\Quote{有一种悔改缺乏宽恕，那么这种悔改也必然完全不能由你们实践。因为无用之事不可行。现在悔改若没有宽恕，便是徒然实践。但\Emph{所有}悔改\Emph{都}应当实践。因此，让我们（承认）\Emph{所有}悔改都获得宽恕，免得徒然实践；因为若徒然实践，便不应当实践。因为若缺乏宽恕，它就是徒然实践。}他们这样指控我们，是公正的；因为他们擅自将这种悔改和其它悔改的果实——即宽恕——占为己有；就\Emph{他们}而言，在其手中（悔改）获得\Emph{人的}平安，（那是徒然的）。然而至于\Emph{我们}，我们记得唯独主赐予罪的宽恕（当然包括\Emph{大罪}），悔改将\Emph{不}会徒然实践。因为（悔改）被交回给主，从此俯伏在主前，正是因这一事实，它更易赢得宽恕，因为它靠恳求\Emph{唯独从天主那里}获得，它不相信\Emph{人的}平安足以抵消其罪债，且就教会而言，它宁愿选择羞耻的红晕，也不愿选择共融的特权。因为它站在教会的门前，以自身耻辱的标记告诫所有人，同时召请弟兄们的眼泪帮助自己，并以更丰富的商品——即他们的同情——而非他们的共融返回。若它在这里没有收获平安的收成，却与主一同播种了平安的种子；它并没有失去，而是准备了它的果实。只要它不疏于本分，它就不会缺少利益。因此，这种悔改并非徒然，这种纪律并非严苛。两者都光荣天主。前者，因不自欺，将更易获得成功；后者，因不自以为是，将更充分地提供帮助。\par

% structure: p0011 source=h2 target=section reason=relative-heading-rank; base=h2

\section{第四章 奸淫与淫乱同义}

既然我们已界定了悔改（种类）的区别，我们现在能回到对罪的评估——看它们是否能在人手之下获得宽恕。首先，（至于事实）我们也将奸淫称为淫乱，习惯要求我们这样做。\Quote{信仰}也与各种名称相熟。所以，在我们每一篇小作品中，我们都谨慎地守护习惯。此外，如果我说\OriginalTerm{奸淫}{adulterium}，以及\OriginalTerm{淫乱}{stuprum}，对血肉玷污的指控将是同一个。因为一个人袭击他人的新娘或寡妇，只要那女子不是他自己的，并无区别；正如地点并无区别——无论贞洁是在内室还是在塔楼中被残杀。每一个杀人，即使发生在树林外，也是强盗行径。因此，凡是享受婚姻之外交合的人，无论在何处，与何种女子，都使自己犯有奸淫和淫乱之罪。相应地，在我们中间，隐秘的结合——即，未曾先在教会面前公开宣称的结合——也有被判断为类似奸淫和淫乱的风险；我们也不得让它们，如果后来被婚姻的覆盖编织在一起，逃脱指控。但是，其他一切情欲的疯狂——对肉体和性别都是不敬的——超出自然规律，我们不仅将它们从门槛驱逐，而且从教会的一切庇护所驱逐，因为它们不是罪，而是异形。\par

% structure: p0013 source=h2 target=section reason=relative-heading-rank; base=h2

\section{第五章 论十诫中禁止奸淫}

那么，奸淫——按其罪行性质，也属于淫乱——究竟该被视为多么深重的罪过，首先有天主的法律来向我们表明；如果确实（事实也是如此），在禁止崇拜异邦神祇和制造偶像本身之后，在提倡（宗教敬礼中）对安息日的尊敬之后，在命令对父母的宗教孝爱仅次于对天主之后，（法律）便紧接着，为了加强和巩固这些诫命，设下了另一条诫命，即\BibleQuote{你不可奸淫。}因为属灵的贞洁与圣德之后，便是身体的完整。（法律）遂立即禁止它的敌人——奸淫，以巩固身体之完整。因此，要明白，那被（法律）规定次于偶像崇拜而予以压制的，是何等样的罪。第二者与第一者并不遥远；没有什么比第二者更接近第一者。由第一者所产生的，在某种意义上便是另一个第一者。所以奸淫与偶像崇拜相邻。因为偶像崇拜也常常以奸淫和淫乱之名被加于百姓身上作为责备，它在命运和跟随上都与奸淫相连——在定罪和安排上同为共犯。不仅如此：（法律）在颁布\BibleQuote{你不可奸淫。}之后，紧接着说\BibleQuote{你不可杀人。}显然，它尊崇了奸淫，使之凌驾于杀人之上，置于至圣法律的最前沿，在天上法令的首要条款中，以最主要之罪的铭刻作记号。从它的位置，你可看出其度量；从其等级，看出其地位；从其邻接，看出其价值。就连邪恶也有其尊荣，即被安置在极恶的顶峰或中央。我看见了奸淫的某种排场：一边，偶像崇拜在前引路；另一边，杀人随行陪伴。无疑，它理所当然地坐在了两种最显著罪行的高峰之间，完全填满了它们之间的空隙，其罪恶的威严同等。被这样的侧翼包围，被这样的肋骨环绕和支持，谁能将它从相互关联的整体、从邻罪的联系、从同类邪恶的怀抱中分离出来，使之独享悔改的益处？难道不正是偶像崇拜在一边，杀人在另一边，将其拘留，并（如果它们有声音）呼喊说：\Quote{这是我们的楔子，这是我们结合的力量？我们以偶像崇拜为标准来衡量；因她的分离介入，我们被联合；向她，从我们中间突出，我们被联结；圣经使我们成为一体；这些文字就是我们的胶水；她自己不能再没有我们而存在。\InnerQuote{我，偶像崇拜，多次为奸淫提供机会；看看我的丛林、我的山丘、活水泉源，以及城市中的圣殿，我们是多么强大的力量，能够推翻端庄。}\InnerQuote{我，杀人，有时也为奸淫效力。且不谈悲剧，看看今日的投毒者，看看魔法师，我报复了多少诱惑，报了多少仇恨；我除掉了多少守卫、多少告密者、多少同谋；再看看接生婆，有多少奸淫的受孕被扼杀。}就算在基督徒中，没有我们，奸淫也不能发生。哪里有不洁之灵的工作，哪里就有偶像崇拜；哪里有人因被污染而被杀害，哪里也有杀人。因此，悔改的治疗方法不适用于\Emph{它们}，否则也适用于\Emph{我们}。我们要么拘留奸淫，要么跟随她。}这些话是罪自身所说的。如果罪缺乏言语，就在（教会门外）附近站着一个偶像崇拜者，附近站着一个杀人者；在他们中间也站着一个奸淫者。同样，如悔改的义务所命，他们穿着苦衣、披着灰烬；以同样的哭泣哀叹；以同样的祈祷巡回；以同样的双膝恳求；呼求同一位母亲。最温良最人道的纪律啊，你在做什么？要么，对\Emph{所有这些}人都应当如此，因为\BibleQuote{缔造和平的人是有福的；}要么，若非对\Emph{所有人}，你就该站在我们这边。你是否一劳永逸地定偶像崇拜者和杀人者的罪，却将奸淫者从他们中间取出？——（奸淫者）是偶像崇拜者的后继者、杀人的先行者、每位的同伙？这是\Quote{偏袒；}你将更可怜的悔改者（无人怜悯地）留在了后面！\par

% structure: p0015 source=h2 target=section reason=relative-heading-rank; base=h2

\section{第六章 旧约时代此类罪行的例子并非新约门徒的典范。但即便旧约中也有对此类罪行的惩罚。}

显然，你若能以天上先例和戒律的庇护，证明你唯独为奸淫——其中也包括淫乱——敞开了悔改之门，那么我们敌对的交锋在此处立即就会兵刃相接。然而，我必然要给你制定一条律法：不可伸手向旧事，不可回头看；因为依照\Person{依撒依亚}，\BibleQuote{旧事已过}；依照\Person{耶肋米亚}，\BibleQuote{更新已更新}；依照宗徒，\BibleQuote{忘记背后，努力向前}；依照主，\BibleQuote{法律及先知至若翰为止}。因为即使我们此刻以法律来证明奸淫之性质，也正合乎基督\BibleQuote{没有废除，而是成全}的那一阶段的法律。因为法律的\BibleQuote{重担}是\BibleQuote{至若翰为止}，而非其疗效之德。被抛弃的是\BibleQuote{行为}的\BibleQuote{轭}，而非纪律之轭。\BibleQuote{在基督内的自由}并未伤害无辜。虔敬、圣德、人道、真理、贞洁、正义、怜悯、仁爱、端庄的法律仍然完整；在这法律中，\BibleQuote{昼夜默思的人是有福的}。关于这法律，同一\Person{达味}又说：\BibleQuote{上主的法律是无瑕的，能转化灵魂；上主的规诫是正直的，能悦乐人心；上主的诫命是光明的，能照亮眼目。}同样，宗徒也说：\BibleQuote{这样，法律是圣的，诫命也是圣的、至善的}——当然是指\BibleQuote{不可奸淫}。但他先前还说过：\BibleQuote{那么，我们因信德废了法律吗？绝对不是！反而坚固了法律}——诚然是在那些方面，因为新约仍禁止它们，且以更严厉的诫命禁止：代替\BibleQuote{不可奸淫}的是\BibleQuote{凡注视妇女而怀贪欲的，已在心里奸淫了她}；代替\BibleQuote{不可杀人}的是\BibleQuote{凡向弟兄说\InnerQuote{拉加}的，应受地狱的罚}。你自问：不可奸淫的法律是否仍然有效？它已加上了不可放纵贪欲的诫命。此外，若有任何（旧约的）先例在你心中偏袒你，它们也不得反对我们所维护的纪律。因为若因古时在某些情况下曾宽恕奸淫，便以为今日也当宽恕，那么附加的法律——定罪罪的\Emph{根源}，即贪欲和意愿，如同实际行为一样——就形同虚设。\Quote{如果较长的先例允许给你的淫荡授予宽免，那么当今更完备的纪律所受的限制又会有什么奖赏呢？}那样，你也会宽恕拜偶像者及一切背教者，因为我们发现百姓本身屡犯这些罪，却又屡次恢复先前的特权。你也会与杀人者共融，因为\Person{阿哈布}藉恳求洗去了\Person{纳波特}的血；\Person{达味}藉忏悔涤除了\Person{乌黎雅}之死及其原因——奸淫。那样，你也会为了\Person{罗特}而宽恕乱伦；为了\Person{犹大}而宽恕结合乱伦的淫乱；为了\Person{欧瑟亚}而宽恕与妓女的卑贱婚姻；为了我们的祖先而宽恕不仅多次结婚，且同时多妻——因为当然，若以古时的先例为理由要求宽恕\Emph{奸淫}，那么对于古时曾获宽恕的一切行为，恩宠也应完全平等。其实，我们也从同一古代拥有支持我们意见的先例——那些不仅没有免除判决，反而立刻对淫乱施以极刑的先例。当然，一个充分的例子是，当百姓与\Place{米德扬}的女儿行淫时，如此众多的人——两万四千人——在一次瘟疫中倒下。但是，为光荣基督，我宁愿从基督引出我的纪律。姑且承认古时——若\OriginalTerm{心性派}{Psychics}喜欢的话——甚至有一切不端庄的（放纵）\Emph{权利}；姑且承认在基督之前，肉身可以嬉戏，甚至可以在它的主去寻找并领回它之前就已\Emph{丧亡}：它那时还不配获得救恩的恩赐，还不适于圣德的职务。直到那时，它仍被视为\Emph{在亚当内}，带着其邪恶的本性，轻易放纵它对所见\BibleQuote{悦人眼目}之物的贪欲，并回头看那低处，用无花果叶搔痒其欲。情欲的毒普遍内在于其中——那由奶形成的渣滓含有它——（这渣滓）适宜如此，因为连水本身也尚未被洗涤。但当天主的圣言降入肉身——（肉身）甚至未曾因婚姻而开启——且\BibleQuote{圣言成了肉身}——（肉身）决不再因婚姻而开启——它要走向的树不是不节制的树，而是坚忍的树；它要从那树尝到的不是甘甜，而是苦涩；它要属于的不是地狱，而是天堂；它要穿戴的不是淫欲的树叶，而是圣德的花朵；它要将自己的纯洁授予水——从那时起，凡\BibleQuote{在基督内}的肉身已失去其原来的污秽，如今成了新造之物，不再由自然种子的污泥或贪欲的污秽所生，而是由\Quote{净水}和\Quote{洁圣的神}所生。因此，你为何要凭古时的先例为它辩护呢？当它曾为奸淫获得宽恕时，它尚未拥有\Quote{基督的身体}、\Quote{基督的肢体}、\Quote{天主的宫殿}这些名号。因此，从它改变了状态，\BibleQuote{受洗归于基督，穿上基督}，且\BibleQuote{以重价赎回}——即\BibleQuote{主和羔羊的血}——的那一刻起，你若抓住任何一个先例（无论是诫命、法律或判决），说宽恕奸淫和淫乱，或过去曾宽恕，或未来将宽恕，那么你也在我们这里得到了一个问题时代的起点定义。\par

% structure: p0017 source=h2 target=section reason=relative-heading-rank; base=h2

\section{第七章 论亡羊与失钱币的比喻}

你可以先从比喻开始，其中主寻找迷失的羊，并把它扛在肩上带回。让你杯上的绘画亲自出来证明，是否在其中那只羊的比喻意义也会透过（外在的表象，用以教导）闪耀出来，表明它所指向的恢复对象是基督徒还是外邦罪人。因为我们根据自然教导、耳目和舌头之律、心智健全而提出异议：这样的回答总是应（问题）而生的，也就是说，是针对提出这些问题的问题而发的。在我看来，引出（答案）的是法利塞人忿忿不平地议论主接纳外邦税吏和罪人与祂来往，并与他们一同进食这一事实。当主用比喻回答这事，设喻描绘迷失之羊的恢复时，我们还能相信祂所指的除了那迷失的\Emph{外邦人}——当时所讨论的就是他们，而不是当时尚不存在的\Emph{基督徒}——还能是谁呢？否则，主岂不成了一个狡辩的回答者，忽略当前本应驳斥的主题，而将精力耗费在将来之事上，这算是什么假设呢？\Quote{但\InnerQuote{羊}严格说来是指基督徒，而主的\InnerQuote{羊群}是教会子民，\InnerQuote{善牧}是基督；因此，在\InnerQuote{羊}身上我们必须理解为一个偏离教会\InnerQuote{羊群}的基督徒。}那样的话，你就让主没有回答法利塞人的议论，而是回应了你的臆测。然而你不得不在辩护这臆测时，否认你认为适用于基督徒的那些（要点）也可以适用于外邦人。告诉我，全人类不都是天主的同一羊群吗？同一个天主不也是万民的主和牧者吗？谁比外邦人更远离天主而\InnerQuote{丧亡}，只要他\InnerQuote{迷途}？当外邦人被基督召回时，谁比外邦人更被天主\InnerQuote{重新寻找}？事实上，这个顺序首先是在外邦人中实现的；也就是说，基督徒的形成无非是先\InnerQuote{迷失}，被天主\InnerQuote{重新寻找}，再由基督\InnerQuote{带回}。同样，这个顺序应当保持，以便我们在解释这类（形象）时，参照那些首先实现它的人。但是，我想你会希望：祂描绘的羊不是从羊群中迷失，而是从方舟或柜子中失落！同样，尽管祂称其余的外邦人为\InnerQuote{义人}，但这并不意味着祂表明他们是\Emph{基督徒}；因为祂是在与\Emph{犹太人}打交道，并在那一刻斥责他们，因为他们对得救的希望感到愤慨。但为了表达，与法利塞人的嫉妒相反，祂自己的恩宠与善意甚至延伸到一个外邦人，祂宁愿一个罪人因悔改而得救，胜过那些凭义德得救的人；否则，请问，难道犹太人\Emph{不}是\InnerQuote{义人}、\InnerQuote{无需悔改}的人吗？因为他们拥有\Quote{\WorkTitle{法律和先知}}作为纪律的舵手和恐惧的工具？因此，祂在比喻中把他们描绘成——即使不像他们实际那样，却像他们应当成为的那样——好让他们在听说悔改对他人必要而对他们不必要时，更加脸红。\par

同样，对于失钱币的比喻，由于是从同一主题引出的，我们也同样解释为指向外邦人；尽管它是在房子里\InnerQuote{遗失}的，仿佛在教会中；尽管它是借着\InnerQuote{灯}\InnerQuote{找到}的，仿佛借着天主的圣言。不，但整个世界是所有人之唯一房屋；在这世界中，更多的是在黑暗中被找到的外邦人，被天主的恩宠光照，而非已经在天主光明中的基督徒。最后，归给羊和钱币的是\InnerQuote{一次}\InnerQuote{迷失}（这对我有利）；因为如果这些比喻是为一个\Emph{基督徒}罪人而设，在他失去信德之后，就应当记载他们的\Emph{第二次}迷失与恢复。\par

我现在暂时从这立场退下；甚至藉着退下，我更能推荐这立场，当我连这样也成功驳倒了对方的狂妄时。我承认每个比喻中描绘的罪人已经是基督徒；但并不是因此就必须肯定，这样的人能藉着悔改从奸淫和淫乱的罪恶中恢复。因为虽然他被称为\\Quote{丧亡}，但还须处理丧亡的\\Emph{种类}；因为\\Quote{母羊}\\Quote{丧亡}不是由于死亡，而是由于迷失；\\Quote{达玛}不是由于毁灭，而是由于隐藏。在这个意义上，一个安全的事物可以被称为\\Quote{丧亡}。因此信徒也\\Quote{丧亡}了，当他从（正路）滑入公开表演的赛马狂热、或角斗士的血腥、或舞台的污秽、或竞技的虚荣；或者若他曾将任何特殊的\\Quote{好奇之艺}用于体育比赛、异教庆典的宴饮、公务需要、或服务于他人的偶像崇拜；若他因某种含混的否认或亵渎之言而自陷。因这类原因，他被逐出羊群；或者他自己，也许因愤怒、骄傲、嫉妒，或——正如常发生的——因不屑接受惩戒，而离开了（羊群）。他应当被寻回并召回。那可以找回的事物并不\\Quote{丧亡}，除非它执意留在外面。藉着在罪人\\Emph{仍活着时}召回他，你将很好地解释这比喻。但是，对于奸淫者和淫乱者，有谁不是在他们犯罪之后立刻宣布他们\\Emph{死亡}？你有什么脸面将那已死的人，依据那召回一只\\Emph{未}死亡的羊的比喻，重新带回羊群？\par

最后，如果你记得先知们，当他们斥责牧者时，有句话——我想是\\Person{厄则克耳}的话：\\BibleQuote{牧者，看哪，你们喝牛奶，穿羊毛；强壮的你们宰杀了，软弱的你们没有牧养；破碎的你们没有包扎；被赶出的你们没有领回；丧亡的你们没有寻找。}请问，他是否也责备他们关于那\\Emph{死亡}的羊，说他们没有注意将那羊也带回羊群？显然，他额外责备他们使羊丧亡，并被田野的野兽吞吃；而它们若留着，既不能\\Quote{丧亡致死}，也不能\\Quote{被吞吃}。\\Quote{难道不可能——（假设）丧亡致死并被吞吃的母羊可以找回——（也依照那在教会——天主的家内达玛（失而复得）的榜样）有些中等性质的罪，相当于达玛的小尺寸和重量，潜伏在同一个教会内，后来被同一个教会发现，便立即在同样伴随的喜乐中结束于悔改？}但奸淫和淫乱不是达玛，而是塔冷通（来衡量）；为搜索它们，需要的不是灯投射的光，而是整个太阳如矛般的光线。这样的人一出现，便被逐出教会；他不在那里存留；他也未给发现他的教会带来喜乐，而是忧伤；他未邀请邻人的祝贺，而是周围兄弟团体的同忧。\par

通过比较，即使以这种方式，将我们的解释与他们的解释相比，母羊和达玛的论点将更指外邦人，以至于它们不可能适用于犯了那罪的基督徒，而他们正是为了这罪而将其强行牵强地应用于对方基督徒身上。\par

% structure: p0023 source=h2 target=section reason=relative-heading-rank; base=h2

\section{第八章 论荡子}

不过，大多数比喻解释者都被同样的结果所欺骗，就如刺绣紫色彩衣时常见的情况一样。当你以为已明智地调和了颜色的比例，相信自己成功巧妙地使它们的相互组合生动起来；而随后，当每块（色彩）和（各种）光完全显现时，那被揭露的差异将暴露所有的错误。因此，在同样的黑暗之中，关于两个儿子的比喻，他们也被其中一些与当前（状态）色调和谐的比喻所引领，偏离了比喻主题所呈现之比较的真正光明之路。因为他们设定两个儿子代表两个民族——长子代表犹太民族，幼子代表基督徒：因为他们无法在后续部分为基督徒罪人（以幼子身份）获得宽恕，除非先以长子身份描绘犹太人。现在，如果我能成功证明犹太人不适合作为长子的比喻，那么结果当然是，基督徒也不能被幼子的共同形象所接纳。因为虽然犹太人也被称为\Quote{儿子}，而且是\Quote{长子}，因为他在义子名分上优先；虽然他也嫉妒基督徒获得天主圣父的和好——这一点对方最热切地抓住——但没有任何犹太人会对圣父说：\BibleQuote{看，我服侍您这么多年，从未违背过您的命令。}因为犹太人何时\Emph{不是}法律的违犯者？听而不闻；憎恨在城门口责备他的人，藐视圣洁的言语？同样，圣父也不会对犹太人说：\BibleQuote{你常与我同在，我的一切都是你的。}因为犹太人被宣告是\BibleQuote{背逆的儿子，确实被生养高升，却不认识主，完全离弃了主，惹怒以色列的圣者。}诚然，我们承认所有事物，显然都是\Emph{赐予}犹太人的；但他同样已从喉咙里被夺走每一块更美味的食物，更不用说父应许之地了。因此，今日的犹太人，不亚于幼子，挥霍了天主的财物，在异乡成为乞丐，直到如今仍服侍那里的王子，即此世的王子。所以，请为基督徒寻找另一个人作为他们的兄弟；比喻不接纳犹太人。他们本可以更恰当地将基督徒比作长子，将犹太人比作幼子，按照信德的类比，如果从丽贝卡腹中所暗示的每个民族的次序允许颠倒的话：只是（那样的话）结尾段落会反对他们；因为基督徒应该为以色列的复兴而喜乐，而不是忧伤，如果我们的全部希望确实与以色列的剩余期望紧密相连，那是真的。因此，即使比喻中某些（特征）有利，但通过其他相反意义的特征，这一比较的彻底贯彻也会被摧毁；尽管（即便所有点都能如镜子般精确对应），解释中仍有一个主要危险——即我们的比较的恰当性被引向不同的目标，偏离了每个特定比喻的主题所命令我们（去调和）的目标。因为我们也记得（看到过）演员，在将寓言姿态与其歌词配合时，表达出与当前情节、场景、角色极不相同的内容，\Emph{却极其协调一致}。但抛开非凡的巧思，因为它与我们的主题无关。同样，异端者也在他们愿意的地方应用同样的比喻，并在不在他们\Emph{应当}的地方——以最大的灵活性——排除它们。为何有最大的灵活性？因为他们从一开始就将他们教义的主题本身按照比喻的偶然巧合塑造在一起。由于脱离了真理法则的约束，他们当然有闲暇去探究和拼凑那些比喻似乎（象征）的事物。\par

% structure: p0025 source=h2 target=section reason=relative-heading-rank; base=h2

\section{第九章 比喻解释的若干一般原则：这些原则应用于当前讨论的比喻，尤其是荡子的比喻。}

然而，我们并不以比喻作为拟订主题的来源，而是以主题作为解释比喻的来源，因此我们也不费心在诠释中（强行）扭曲一切（使之成形），同时留意避免一切矛盾。为什么\Quote{一百只羊？}又为什么，自然是\Quote{十块德拉克马？}那\Quote{扫帚？}又是什么？那愿意表达一个罪人的救恩给天主带来的极大喜乐者，必须说出某个整体数量中的特定数目，以描述那\Quote{一}已经丧亡。在\Quote{屋内}寻找\Quote{德拉克马}的人，其过程必须适当配合\Quote{扫帚}和\Quote{灯}的有益辅助。因为这类好奇的细节不仅使一些事受质疑，而且因牵强解释的狡黠，常导致偏离真理。此外，有些要点仅仅是为了比喻的结构、安排和纹理而引入，以便它们能被通篇加工，以达到提供典型例子的目的。当然，两个儿子的比喻与德拉克马和母羊的比喻指向同一结局：因为它与它们所结合的比喻有相同的原因（引起它），当然也有法利塞人因主与外邦人来往而发出的同样的\Quote{嘀咕}。否则，若有人怀疑在\Place{犹太地}（久已被\Person{庞培}和\Person{卢库鲁斯}之手征服）税吏是外邦人，就让他读\BibleBook{}{申命纪}：\BibleQuote{以色列子民中不应有税吏。}税吏这名号在主眼中也不会如此可憎，除非它是一个\Quote{陌生}的名号——即那些将天空、大地和海洋的道路都拿来出卖之人的名号。此外，当（作者）将\Quote{罪人}与\Quote{税吏}并列时，并不表示他显示他们是犹太人，尽管有些可能如此；但他通过将外邦人的同一\Emph{类属}——有些是因职务而成为罪人，即税吏；有些是因本性而成为罪人，即非税吏——置于同等地位，而对他们做出了区分。此外，主不会因与犹太人一同吃饭而受责备，而是因与外邦人一同吃饭——犹太人的纪律将（其门徒）排除于外邦人的餐桌之外。\par

现在我们必须继续论述荡子的比喻，首先考虑那更有益的部分；因为若以最精细的天平调整比喻，却证明对救恩极为有害，便不得认可。但我们看到，整个救恩体系，正如它包含在纪律的维持中，正被反对者所偏爱的解释所颠覆。因为倘若一位\Emph{基督徒}远离他的圣父，过着外邦人的生活，挥霍了从天主圣父所领受的\Quote{实质}——当然是圣洗的实质，圣神的实质，以及由此而来的永生的希望——倘若他被剥夺了精神的\Quote{财产}，甚至将自己的服事交给了世界之王——除了魔鬼还能是谁？——并且被他指派从事\Quote{放猪}的事务，即照料不洁的邪灵，以致他恢复理智回到圣父那里——其结果将是：不仅是奸淫者，而是偶像崇拜者、亵渎者、叛教者以及各类背教者，都将借此比喻向圣父作补赎；这样，或许可以说，整个圣事的\Quote{实质}就最真实地被浪费了。因为谁会害怕挥霍那些以后还能再恢复的东西呢？谁会小心翼翼地永远保存那些\Emph{并非}永远可能失去的东西呢？罪中的安全同样是对罪的欲望。因此，背教者也将恢复他先前的\Quote{外衣}，即圣神的袍服；重新得到\Quote{戒指}，即圣洗的标记和印记；基督将再次被\Quote{宰杀}；他将躺卧在那张床榻上，而凡是\Emph{不相称地穿着衣服}的人，往往被刑役提起并抛入黑暗中——更何况那些被\Emph{剥去}衣服的人呢。因此，更进一步说，如果荡子的故事不适用于基督徒，那它既不\Emph{适宜}，也不\Emph{合理}。所以，即使\Quote{儿子}的形象也不完全适用于犹太人，我们的解释只需以主的目的为准。主来当然是为了拯救那\Quote{丧亡的}，他是\Quote{医生}，为\Quote{有病的人}比\Quote{健全的人}更为必要。这事他惯常既以比喻预表，又以直言宣讲。人中有谁\Quote{丧亡}，谁从健康中坠落，除非是不认识主的人？谁是\Quote{安全无恙}的，除非是认识主的人？这两类人——生来即为\Quote{兄弟}——这个比喻也将指出来。请看外邦人是否在天主圣父那里拥有起源的\Quote{实质}、智慧以及认识天主的天赋能力；宗徒也藉此能力指出，\BibleQuote{在天主的智慧中，世界凭智慧不认识天主}——这智慧当然起初是从天主领受的。因此，他挥霍了这\Quote{实质}；因他的道德习性远离了主，陷入世界的错误、诱惑和欲望中，在那里，因对真理的饥渴所迫，他将自己交给了世界之王。世界之王派他看管\Quote{猪}，放牧那与邪灵相熟的畜群，在那里他无法掌握生命之粮的供给，同时却看到他人从事神圣工作，拥有充足的天上食粮。他想起他的圣父，天主；当他满足后，便回到他那里；重新得到那原始的\Quote{外衣}——即亚当因过犯所失落的状态。他也往往在那时初次得到\Quote{戒指}，用以在受问时公开封印信仰的盟约，从此以后就享有主身体的\Quote{肥甘}——即圣体圣事。这就是荡子，他以前从未节俭；他从一开始就是放荡的，因为他从起初就\Emph{不是}基督徒。当他从世界回到圣父的怀抱时，法利塞人藉着\Quote{税吏和罪人}的身份为他忧伤。因此，长兄的嫉妒只适用于这一点：不是因为犹太人无辜并顺从天主，而是因为他们嫉妒外邦人的得救；他们显然正是那些本应\Quote{常与}圣父同在的人。当然，犹太人是在基督徒的\Emph{第一次}蒙召时悲叹，而不是在基督徒的\Emph{第二次}复原时；因为前者甚至照耀外邦人；而后者发生在教会中，甚至连犹太人也不知道。我认为我提出了更符合比喻主旨、事物一致性以及纪律保持的解释。但如果对方急切地将羊、银币和儿子的纵情享乐塑造成基督徒罪人的样式，是为了赋予奸淫和私通以悔改的机会；那么，要么所有其他同样严重的罪都应被允许赦免，要么与它们同类的奸淫和私通应被保留而不予许可。\par

但更重要的是：将结论推广到与当前直接讨论的主题无关的事情上是不允许的。简言之，如果允许将比喻转用于（原意之外的）其他目的，那么我们要将这些比喻所唤起的希望指向\Emph{殉道}更为合理；因为唯有殉道，在浪子耗尽一切家产后，能使他重获儿子名分；会喜乐地宣告那德拉克马已经找到，即便它曾混在粪堆的垃圾中；会将那尽管在崎岖之地四处流亡的母羊，背在\Person{主}自己的肩上带回羊群。但是，如果必须如此选择，我们宁愿在圣经上\Emph{少}一点智慧，也不愿\Emph{悖逆}圣经而自诩智慧。我们有责任保持\Emph{主的意义}，如同保持祂的\Emph{诫命}。在解释上的逾越，并不比在言语上的更轻。\par

% structure: p0029 source=h2 target=section reason=relative-heading-rank; base=h2

\section{第十章  悔改更适用于外邦人而非基督徒}

因此，当那禁止将这些比喻用于外邦人的轭被挣脱，并且终于看出或承认必须按照命题的主题进行解释之后；他们接着争辩说，悔改的正式宣告甚至不适用于外邦人，因为他们的罪不适用于悔改，这些罪源自无知，仅是本性就使他们在天主面前成为有罪的。因此，对于那些连危险本身都不了解的人，补救措施也是不可理解的；然而，悔改的原则存在于罪是带着良心和意愿所犯的地方，那里过错和恩宠都是可理解的；那哀恸、那俯伏的人，他知道他失去了什么，以及如果他把自己的悔改作为祭品献给天主——祂当然是向儿子而非陌生人提供悔改——他将恢复什么。\par

那么，这就是\Person{约纳}在拖延传道职责时，认为悔改对尼尼微的外邦人不必要的理由吗？或者他倒是预见到天主的仁慈甚至倾注在陌生人身上，担心那仁慈会（可以说）毁坏他宣告的（信誉）？而为了一个尚未认识天主、仍在无知中犯罪的世俗城市，先知几乎丧命？\EditorialNote{约 1:4}，只是他预尝了\Person{主}苦难的典型预像，这苦难是要在他们悔改时赎回外邦人（如同其他人）。对我而言，就连\Person{若翰}在\Quote{铺平主的道路}时，也向服兵役的人和税吏，如同向亚巴郎的子孙一样，宣告悔改，这已足够。\Person{主}自己也假定，西顿人和提洛人若看见了祂\Quote{奇迹}的证据，就会悔改。\par

不，我甚至要主张，悔改更适用于本性的罪人，而非故意的罪人。因为尚未\Emph{使用}过悔改的人，比已经\Emph{滥用}过它的人，更配得它的果实；第一次使用的补救措施比用旧了的更有效。毫无疑问，\Person{主}是\Quote{仁慈的}对待\Quote{忘恩的人}，而不是对待无知的人！祂\Quote{慈悲地}对待\Quote{被遗弃的人}更甚于对待尚未被试探的人！以至于对祂仁慈的侮辱，难道不会招致祂的\Emph{愤怒}，而非祂的\Emph{抚慰}！祂岂不更愿意将那（仁慈）赐给陌生人，而这仁慈在祂自己儿子们身上，祂已经失去，因为犹太人嘲弄祂的忍耐，祂就这样收养了外邦人！但\Quote{属血气的人}的意思是说：天主，公义的审判者，更愿罪人悔改而非死亡，那罪人却选择了死亡而非悔改！若真是这样，我们便是借犯罪来赢取恩宠了。\par

来吧，你这在端庄、贞洁和每一种性的圣洁上走钢丝的人，借着这种远离真理之途的规诫之助，用不稳的脚步攀登在最纤细的线上，以肉身平衡灵性，以信德节制你的兽性原则，以敬畏调适你的眼睛；为何你如此全神贯注于一步？继续前进，如果你成功找到了能力和意志，而你是如此安稳，仿佛立于坚实之地。因为若有任何肉身的摇摆、心神的分散或眼目的游荡，将你从平衡中震落，\Quote{天主是良善的。}祂不是向外邦人，而是向自己的（子女）敞开怀抱：第二次补赎将等待你；你将再次从一个奸淫者成为基督徒！你（会）将这些（恳求）用在我身上，你这最和蔼的天主诠释者。但若那唯一偏袒奸淫者的经书\WorkTitle{牧者}值得在神圣正典中占有一席之地，我就向你退让；但它岂非被每一届教会（甚至你们自己的）会议惯常判定为伪经和虚假（文献）？它本身是奸淫的，因而成为其同伙的保护者；你在其他方面也从中获得启蒙；也许，那个\WorkTitle{牧者}会成为庇护者，就是你描绘在你（圣事性）圣爵上的那位，（我说，描绘他）他自己也是基督教圣事的卖淫者，（因此）配得既是醉酒的偶像，也是奸淫的牛虻，圣爵将迅速随之而来，（这圣爵）你从其中最轻易啜饮到的莫过于第二次补赎的\Quote{母羊}！然而，我畅饮那不能被打破的牧者的圣经。\Person{若望}立刻将祂连同洗礼盆和补赎的责任一并呈给我；（且将祂呈为）说：\BibleQuote{结出与补赎相称的果实，不要心里说：我们有\Person{亚巴郎}为父}——以免他们再次从赐予祖先的恩宠中为过犯沾沾自喜——\BibleQuote{因为天主能从这些石头中给\Person{亚巴郎}兴起子孙来。}因此，我们也须判断那些\BibleQuote{不再犯罪}的人，如同\BibleQuote{结出与补赎相称的果实}。因为还有什么比改正的成就更能成熟为补赎的果实呢？但即使\Emph{宽恕}更是补赎的果实，宽恕也不能在不停止犯罪的情况下共存。同样，停止犯罪是宽恕的根源，以使宽恕成为补赎的果实。\par

% structure: p0034 source=h2 target=section reason=relative-heading-rank; base=h2

\section{第十一章　从比喻开始，\Person{德尔图良}转而考虑主明确的行事。}

从其与福音的关联性来说，比喻的问题此时确实已经处理完毕。然而，如果主凭祂的作为也发布了任何有利于罪人的宣告；例如当祂允许那\Quote{一个女人，是个罪妇}接触祂自己的身体——她用眼泪洗祂的脚，用头发擦干，并用香膏预备祂的安葬；又如当祂向那撒玛黎雅妇人——她并非因第六次婚姻而成为奸妇，而是一个娼妓——（轻易地）显明了祂是谁；（那么）我们的对手并未从中获得任何好处，即使祂（在这些情况下）所赐予宽恕的人已经是基督徒了。因为我们现下断言：这唯独是主所允许的：愿祂宽仁的权能在今日仍然奏效！然而，在祂生活于地上的那些时期，我们明确设定：如果宽恕曾赐予罪人——甚至是犹太罪人——这并不构成不利于我们的预判。因为基督徒的纪律始于盟约的更新，并且（如我们已预先说明的）始于肉身的救赎——即主的受难。在信德秩序被发现之前，无一人是完美的；在基督升天之前，无一人是基督徒；在圣神从天上显现之前，无一人是圣洁的——祂是纪律本身的制定者。\par

% structure: p0036 source=h2 target=section reason=relative-heading-rank; base=h2

\section{第十二章　论宗徒在会议中关于奸淫的裁决。}

因此，那些藉着宗徒并透过宗徒接受了\Quote{另一位\OriginalTerm{护慰者}{Paraclete}}的人——他们既不曾在祂特选的先知身上认出这位护慰者，如今在宗徒身上也不再拥有祂——那么，就让他们现在，甚至从宗徒的著作中，教导我们：那在领洗后重新玷污的肉身，其污秽能否藉着补赎得以洗净。难道我们不是在宗徒身上也认出了旧法律关于指出奸淫之重大罪过的形式吗？免得它在新的训诫阶段中被视为比旧阶段轻微。当初福音如雷霆般震动、将旧制度摇动至根基之时，当争论法律是否应被保留之际，宗徒们藉着圣神的权威，向那些已开始从外邦人中聚集到他们身边的人发出第一条规则：\BibleQuote{圣神和我们决定，不加给你们任何更重的担子，只限于那些必要遵守之事：即祭偶像的物、淫乱和血；你们戒绝这些，就做得好，圣神带领着你们。}在此处，奸淫和淫乱已足够被保留在崇拜偶像与凶杀之间那属于其自身的尊贵地位：因为关于\Quote{血}的禁令，我们更应理解为是对\Emph{人的}血的禁令。那么，宗徒们是如何希望那些他们从原初法律中特意挑选出来、要小心防范的罪行呈现的呢？他们规定哪些是必须禁绝的呢？并非他们容许其他罪行，而是他们唯独将这些罪行置于最前列，当然作为不可赦免的；他们为了外邦人的缘故，使法律的其他负担成为可赦免的。那么，他们为什么从如此沉重的轭下释放我们的颈项，除非是为了将这些训诫的纲要永远加在这些颈项上？他们为什么宽容地放松如此多的束缚，除非是为了使我们永久地完全受制于那些更必要的？他们从较多的束缚中释放了我们，好使我们能被束缚起来，禁绝那更有害的。此事已通过补偿而确定：我们得到了许多，好使我们能付出一些。但这补偿是不可撤销的；也就是说，如果它被重复——当然是奸淫、血和崇拜偶像的重复——所撤销：因为如果宽恕的条件被违反，那么整个法律的负担就会被重新承担。但圣神与我们达成协议，并非轻率——祂甚至在我们没有请求的情况下就达成了这协议；因此祂更应受尊敬。只有忘恩负义的人才会解除祂的约定。在这种情况下，祂既不会收回祂已丢弃的，也不会丢弃祂所保留的。新盟约的条件永远不变；当然，那法令及其所包含的劝谕的公开宣读，将只与世界一同止息。祂已足够明确地拒绝宽恕那些祂特意命令要小心避免的罪行；凡是祂没有藉推论予以许可的，祂都已然要求。因此，教会绝不将平安的恢复赐予\Quote{崇拜偶像}或\Quote{血}。我认为，相信宗徒们会偏离他们自己的最终决定，是不合法的；否则，如果有人觉得可以这样相信，他们就必须证明它。\par

% structure: p0038 source=h2 target=section reason=relative-heading-rank; base=h2

\section{第十三章。论圣保禄，以及他劝格林多人宽恕的那个人。}

在此我们也清楚知道他们所提出的猜疑。事实上，他们怀疑宗徒保禄在\WorkTitle{格林多后书}中，竟宽恕了同一个犯奸淫的人，就是他在第一封信中公开判决要\BibleQuote{交给\Person{撒殚}，为\OriginalTerm{毁灭他的肉身}{destruction of the flesh}}的那个\BibleRef{格前 5:5}——那人是其父亲婚床的大逆不道的继承人；好像后来他抹去了自己的话，写道：\BibleQuote{但若有人使你们完全忧愁，他并没有使我完全忧愁，只是部分忧愁，免得我加重你们的负担。由多数人给予的这种责斥就够了；故此，相反地，你们更应该原谅并安慰他，免得这样的人因过度忧愁而被吞没。为此，我求你们向他们重申\OriginalTerm{友爱}{affection}。我写信给你们，也正是为此，为要知道你们在一切事上是否顺从。如果你们宽恕谁，我也宽恕；因为我若宽恕什么，也是在基督面前宽恕的，免得我们被\Person{撒殚}胜过，因为我们不是不知道他的阴谋。}\BibleRef{格后 2:5-11}这里所指的与那犯奸淫的人有什么关系？与那玷污父亲床铺的人有什么关系？与那越过了外邦人的无耻的基督徒有什么关系？——他当然会用特别的宽恕来赦免一个他曾以特别愤怒谴责的人。他在怜悯中比在义怒中更隐晦，在严厉中比在温和中更公开。然而，一般来说，愤怒比宽容更容易转弯。忧伤的事物比喜乐的事物更容易犹豫。当然，这里所涉及的是某种\Emph{适度}的宽容；这种（宽容中的适度）此时，如果曾有的话，需要被推测出来，因为所有\Emph{最大}的宽容通常都不会在没有公开宣告的情况下被授予，更不用说没有具体指明了。为什么？你自己，当为了用他的祈祷感动弟兄们而将悔改的通奸者带进教会时，你把他领到中间，使他俯伏，穿着苦衣和灰烬，一副羞辱与恐怖的混合物，在寡妇面前、在长老面前，恳求众人的眼泪，舔众人的脚印，抱众人的膝盖？而你，这位好牧人和蒙福的父亲，为了实现这个人的（期望的）结局，用你能力范围内所有怜悯的诱饵来美化你的演讲，并在\Quote{母羊}的比喻下去找寻你的山羊？你，害怕你的\Quote{母羊}再次从羊群中跳出去——仿佛未来的事不再合法，而它甚至连一次都不曾合法——难道在赐予宽恕的那一刻，也使其余所有的人都充满恐惧？而宗徒会如此轻率地宽恕那背负着乱伦的凶残纵欲，甚至没有要求罪犯穿上你所应当从他那里学到的这法定忏悔服？甚至没有对过去发出警告？对未来没有训诫？不仅如此，他更进一步，恳求他们\BibleQuote{向他重申友爱}，仿佛他在赔偿他，而不是在宽恕！然而我听到（他说）\BibleQuote{友爱}，而不是\BibleQuote{\OriginalTerm{共融}{communion}}；正如他在\WorkTitle{得撒洛尼后书}中写道：\BibleQuote{若有人不听从我们这书信上的话，你们要记下他，不要与他交往，使他自愧；但不要以他为仇人，而要劝戒他如同弟兄。}\BibleRef{得后 3:14-15}因此，他也可能对犯奸淫的人说，只允许\BibleQuote{友爱}，而不允许\BibleQuote{共融}；然而对于乱伦的人，连\BibleQuote{友爱}也没有；他必会吩咐人把他们从他们的\Emph{中}驱逐——当然更从他们的\Emph{心}中驱逐。\BibleQuote{但他担心他们因失去那已被他交给\Person{撒殚}的人而被\Person{撒殚}胜过，或者怕那被他判决\InnerQuote{毁灭肉身}的人因\InnerQuote{过度忧伤}而被吞没。}\BibleRef{格后 2:5-11}他们甚至这样解释\Quote{肉身的毁灭}为忏悔的工作；因为通过斋戒、污秽以及各种刻意的忽视和精心虐待，致力于消灭肉身，似乎以此向天主赎罪；所以他们争辩说，那个犯奸淫的人——更确切地说，那个乱伦的人——由宗徒交付给\Person{撒殚}，不是出于\Quote{\OriginalTerm{丧亡}{perdition}}的目的，而是出于\Quote{\OriginalTerm{改正}{emendation}}的目的，基于假设后来他因\Quote{肉身}的\Quote{毁灭}（即普遍的苦难）而得宽恕，因此确实得到了宽恕。显然，同一宗徒把\Person{依默纳约}和\Person{阿肋桑德}交付给\Person{撒殚}，\BibleQuote{使他们受管教不再\OriginalTerm{亵渎}{blasphemy}}，正如他在\WorkTitle{弟茂德前书}中所说。\BibleRef{弟前 1:20}\BibleQuote{但他自己也说：\InnerQuote{有一根刺加给了我，就是\Person{撒殚}的使者}来击打我，免得我过于高举自己。}\BibleRef{格后 12:7}如果他们同样触及这个（例子），为了使我们明白那些被他\Quote{交付给\Person{撒殚}}的人（是这样交付的）是为了改正，而不是为了丧亡；那么，亵渎和乱伦之间有什么相似之处呢？与一个完全免于这些——不，反而因最高圣德和全然无罪而自高自大的灵魂（灵魂的骄矜）有什么相似呢？这种（灵魂的骄矜）在宗徒身上正被\Quote{击打}所抑制，如果你愿意，是通过（如他们所说）耳朵或头部的疼痛。然而，乱伦和亵渎应当将整个人交付给\Person{撒殚}本人作为占有物，而不是交付给\Person{撒殚}的\Quote{使者}。还有（另一点）：关于这一点是有区别的，不，更确切地说，在这一点上至关重要：我们发现那些人是由宗徒交付给\Person{撒殚}，但给宗徒自己却赐了一个\Person{撒殚}的使者。最后，当\Person{保禄}祈求主除去它时，他听到了什么？\BibleQuote{我的恩宠够你用，因为德能在软弱中才显全。}\BibleRef{格后 12:9}那些被交付给\Person{撒殚}的人听不到这个。此外，如果\Person{依默纳约}和\Person{阿肋桑德}的罪行——即亵渎——在今世和来世都是不可赦免的，那么宗徒当然不会违背主的决定，在\Emph{宽恕的希望下}把那些已经从信仰陷入亵渎的人交给\Person{撒殚}；因此他也称他们为\BibleQuote{在信仰上遭了船破}，不再有船只——教会——的慰藉。因为对于那些信了后又触礁（陷入）亵渎的人，宽恕被拒绝；另一方面，\Emph{外邦人}和\Emph{异端者}却每天从亵渎中出来。但即使他说：\BibleQuote{我把他们交给\Person{撒殚}，使他们受管教不再亵渎}，他是对其他人说的，这些人通过被交付给\Person{撒殚}——即被逐出教会之外——必须受训练，知道不可亵渎。同样，那乱伦的犯奸淫者，他也交付给\Person{撒殚}，不是为了改正，而是为了丧亡，因为他已经通过比外邦人更严重的罪，归向了\Person{撒殚}；使他们知道不可犯奸淫。最后，他说：\BibleQuote{为\Emph{毁灭}肉身}，不是\Quote{\Emph{折磨}肉身}——谴责那造成他（从信仰）堕落的具体实体，这实体在丧失圣洗时已经立即丧亡——\BibleQuote{为使灵魂}，他说，\BibleQuote{在主的日子可以得救。}\BibleRef{格前 5:5}（这里又有一个困难）：因为要探究这一点，是否\Emph{那人自己的灵魂}将会得救。如果是这样，一个被如此深重的邪恶所污染的灵魂将会得救；肉身的毁灭是为了使灵魂可以在\Emph{惩罚中}得救。若是这样，与我们相反的解释将承认一个\Emph{没有肉身}的惩罚，如果我们失去了肉身的复活。因此，他的意思仍然是：那被认为存在于\Emph{教会内}的灵魂，必须通过驱逐那乱伦的犯奸淫者，而被呈献为\BibleQuote{得救}，即在主的日子里不受不洁的传染；如果，也就是说，他接着说：\BibleQuote{你们不知道，一点面酵能使整个面团发起来吗？}\BibleRef{格前 5:6}然而乱伦的奸淫不是一点点面酵，而是大块面酵。\par

% structure: p0040 source=h2 target=section reason=relative-heading-rank; base=h2

\section{第十四章　续论同一主题。}

And— these intervening points having accordingly been got rid of— I return to the second of \WorkTitle{格林多后书}; in order to prove that this saying also of the apostle, \BibleQuote{对于这样的人，由许多人施行的\Emph{这次责备}已经足够了}，并不适用于那行淫乱者。因为如果宗徒曾判决他\BibleQuote{要把他交给撒殚，以败坏他的肉体}，那么他当然是\Emph{定罪}了他，而非\Emph{责备}了他。因此，另有一个人是宗徒愿意以\BibleQuote{责备}为足够的；如果那行淫乱者从判决中获得的不是\BibleQuote{责备}，而是\BibleQuote{定罪}。因为我将这个问题交给你考察：在第一封书信中，是否也有其他人因\BibleQuote{行为不端}而\BibleQuote{使宗徒全然忧愁}，并因招致（宗徒的）\BibleQuote{责备}而\BibleQuote{被他全然忧愁}，如同第二封书信所暗示的？其中某一位可能在第二封书信中得到了宽恕。此外，让我们注意整封第一封书信，它（可以这样说）是整体写成的，不是用墨水，而是用胆汁；充满怒气、愤慨、轻蔑、威胁、敌意，并且针对某些个体，形同对他们的具体指控而塑造。因为分党、嫉妒、争论、傲慢、骄傲、纷争要求他们被敌意压制，被简短的责备驳回，被傲慢挫伤，被严厉吓阻。而谦卑的锐利又是什么样的敌意呢？\BibleQuote{我感谢天主，除了\Person{克黎斯颇}和\Person{加约}外，我没有给你们任何人付洗，免得有人说我因自己的名字付洗。} \BibleQuote{因为我曾决定，在你们中间不知道别的，只知道耶稣基督，这被钉在十字架上的基督。} 以及\BibleQuote{（我想）天主把我们宗徒列在最后，好像定了死罪的人，因为我们成了一台戏，给世界、天使和世人观看。} 以及\BibleQuote{我们成了世上的垃圾，众人的废物。} 以及\BibleQuote{我不是自由的吗？我不是宗徒吗？我没有见过我们的主耶稣基督吗？} 反之，他无奈地宣告：\BibleQuote{但对我来说，被你们审断，或被人间的法庭审断，是极小的事；因为我连自己也不判断。} 以及\BibleQuote{我的夸耀没有人能使之落空。} \BibleQuote{你们不知道我们要审判天使吗？} 此外，那自由的表达是何等明显的指责，属灵之剑的锋芒是何等显著：（如这些词句：）\BibleQuote{你们已经饱足了！已经富足了！已经为王了！} 以及\BibleQuote{若有人自以为知道什么，他还不知该怎样知道。} 他岂非甚至\BibleQuote{击打某人的脸}，说道\BibleQuote{是谁使你不同呢？你有什么不是领受的呢？你为什么夸耀，好像不是领受的呢？} 他岂非也在\BibleQuote{击打他们的口}，说道：\BibleQuote{但有些人仍凭着良心，直到如今还吃祭偶像之物。但这样犯罪，使软弱的弟兄良心受到震动，他们便得罪了基督。} 至此，他点名道姓：\BibleQuote{难道我们没有权柄吃喝，并带领一位姊妹往来，如同其他的宗徒、主的弟兄和\Person{刻法}一样吗？} 以及\BibleQuote{如果别人在权柄上有份于你们，难道我们不是更有份吗？} 同样，他用个性化的笔触刺伤他们：\BibleQuote{所以，那自以为站得稳的，要小心，免得跌倒。} 以及\BibleQuote{若有人似乎好争论，我们没有这样的习俗，主的教会也没有。} 以这样的结语，加上诅咒：\BibleQuote{若有人不爱主耶稣基督，就让他当受诅咒。主来！} 他当然是在击打某个特定的人。\par

但我宁愿站在宗徒更为热忱之处，即那行淫乱者本人也使他人困扰之处。\BibleQuote{似乎我不去到你们那里，有些人就自高自大。但主若愿意，我必快来到，并且要知道那些自高自大的人的言语，而是他们的能力。因为天主的国不在于言语，而在于能力。你们愿意怎样？要我带着棍棒到你们那里去，还是怀着慈爱和温柔的精神？} 接着是什么？\BibleQuote{我确实听说你们中间有淫乱的事，这样的淫乱连外邦人中也没有，就是有人有他父亲的妻子。你们还自高自大，而不哀痛，好把行这事的人从你们中间除去？} \Emph{他们应为谁} \BibleQuote{哀痛}？当然是为一个死人。\Emph{向谁}哀痛？当然是向主，好使他以某种方式\BibleQuote{从他们中间被除去}；当然不是被逐出教会。因为凡属于教会首领职分的事，不会向天主祈求；但所求的是，通过死亡——不仅是众人共同的死亡，而是特别适用于那已经是尸体、患有不治之症的坟墓般的肉体——他更能（比单纯绝罚）承受\BibleQuote{从教会中被除去}的惩罚。因此，既然当时他可能被\BibleQuote{除去}，宗徒\BibleQuote{判定这样的人要交给撒殚，以败坏他的肉体}。因为被抛给魔鬼的肉体理应被诅咒，好使之从祝福的圣事中被丢弃，永不再回到教会的营中。\par

因此，我们在此看到宗徒的严厉被区分：一面针对那\Quote{自高自大的}人，一面针对那\Quote{乱伦的}人；（我们看到宗徒）以\Quote{棍杖}武装对付前者，以判决对付后者——一根他正在威胁的\Quote{棍杖}，一项他正在执行的判决：前者（我们看到）仍在挥舞，后者瞬间投掷；（一者）他用之责斥，另一者他用之定罪。确实，紧接着，那受责斥者就在高举的棍杖的威吓下战栗，但被定罪者却在刑罚的即刻施行下灭亡。前者立即退却，畏惧打击；后者则偿付了刑罚。当同一位宗徒的第二封信送达格林多人时，明确地给予了宽恕；但不确定\Emph{给谁}，因为既未指明人，也未指明原因。我将用感官比较这些案件。如果那\Quote{乱伦的}人摆在我们面前，那\Quote{自高自大的}人也会站在同一平台上。当然，案件的类比足够成立，因为\Quote{自高自大的}受责斥，而\Quote{乱伦的}被定罪。对\Quote{自高自大的}给予了宽恕，但在责斥之后；对\Quote{乱伦的}似乎未给予宽恕，因其处于定罪之下。如果宽恕是要赐给那使人担心会被\Quote{被忧伤吞噬}的人，那么那\Quote{受责斥的}人仍处于被吞噬的危险中，因威胁而灰心，因责斥而忧伤。然而，那\Quote{被定罪}的人则被永久视为已被吞噬，既因他的过错也因他的判决；（也就是说，被视为）一个不需要\Quote{忧伤}，而是要\Emph{忍受}那在忍受之前他本可忧伤之事的人。如果给予宽恕的原因是\Quote{以免我们被撒殚欺骗}，那么所预防的损失涉及尚未灭亡之物。对于已经了结之物，无需预防；但对于仍然安全之物，则需预防。但那位被定罪者——也被定罪交付给撒殚——当他在犯下如此行为的那一刻，已经\Emph{从教会中}灭亡了，更不用说在教会本身弃绝他的那一刻。教会怎会害怕遭受他的欺诈性损失呢？她早已在他被夺去时失去了他，并且在定罪之后，她已无法留住他。最后，法官给予宽恕适合针对什么？是针对他通过正式宣告已最终定案之事，还是针对他通过中间判决悬而未决之事？当然，（我指的是）那位不习惯\Quote{重建他所拆毁的，以免被视为犯法者}的法官。\par

来吧，如果他在第一封信中没有\Quote{完全忧苦}那么多人；如果他未曾\Quote{责斥}任何人，未曾\Quote{惊吓}任何人；如果他只\Quote{击打}了那乱伦的人；如果为了他的缘故，他没有使人陷入恐慌，没有使（任何）\Quote{自高自大者}惊惶失措——那么，怀疑有某位完全不同的人当时在格林多人中处于同样困境，岂不是更好吗？而推论他因此——因其过错的适度性质——随后获得了宽恕，岂不是更令人相信吗？这样，你就不必将那宽恕解释为赐给乱伦的邪淫者了。因为即使不是在书信中，你也本应读到这些，它更清楚地印在宗徒的品格上，藉着\OriginalTerm{端庄}{modesty}，而非藉着笔的媒介：也就是说，不要让\Person{保禄}——\Quote{基督的宗徒}，\Quote{在信德和真理中教导万民者}，\Quote{蒙选之器}，教会的建立者、纪律的监察者——陷入如此大的\OriginalTerm{轻率}{levity}，以致他要么草率地定罪了一个他即将赦免的人，要么草率地赦免了一个他并未草率定罪的人，尽管是基于那种源于单纯不端庄的邪淫，更不用说基于乱伦的婚姻、不敬的纵欲和弑亲的贪欲了——他甚至拒绝将这种贪欲与外邦人的（贪欲）相比较，以免被归咎于习俗；他虽不在场，却要审判这种贪欲，以免罪犯\Quote{赢得时候}；他召集了\Quote{主的德能}来辅助自己才定罪，以免判决看似人为的。因此，他既戏弄了自己的\Quote{心神}，也戏弄了\Quote{教会的天使}和\Quote{主的德能}，如果他撤销了在他们的参议下所正式宣布的判决。\par

% structure: p0045 source=h2 target=section reason=relative-heading-rank; base=h2

\section{第十五章　接续同一主题。}

如果你悉心推敲那封书信的续文，以阐发宗徒的意旨，你会发现那段续文与赦免乱伦之举不符；免得在此宗徒也因他后来含义的不一致而蒙羞。因为这是何等的一种假设，即在将教会和平的特权慷慨地恢复给一个乱伦的淫乱者之后，他立刻转而累积起关于远离不洁、除去污点、劝勉圣善行为的劝诫，仿佛他先前并未颁布任何相反性质的法令似的？简而言之，请比较（并看看）是否他有资格说：\BibleQuote{因此，我们既然蒙受了怜悯，拥有了这职分，就不沮丧；反倒弃绝了可耻的隐密事}，这人刚刚释放了一个明显被定罪的人，此人不仅犯了\Quote{可耻}的事，而且犯了罪：同样，他是否有资格去宽恕一种明显的放纵无度，他竟在列举自己的劳苦时，先说到\Quote{困苦和压力}，又说到\Quote{禁食和守夜}，随后也提到了\Quote{贞洁}：再者，他是否有资格去接纳任何邪恶的人回到共融之中，他写着：\BibleQuote{义与不义有什么相交？光明与黑暗又有什么相通？基督与\OriginalTerm{贝里雅耳}{Belial}有什么和谐？信者与不信者有什么相干？天主的圣殿与偶像有什么一致？}他难道不应当不断听到（这样的回答）吗：\Quote{你是用什么方式把你在这封书信前部因恢复乱伦者的共融而联结起来的事物分开的呢？因为因他重新与教会成为一体，义就与不义相交，光明与黑暗相通，贝里雅耳与基督和谐，信者与不信者共享圣事。至于偶像，它们自己看着办吧：连天主圣殿的玷污者也被转化成了天主的圣殿：因为在这里，他也说：\InnerQuote{因为你们是永生天主的殿。正如祂说：我要在你们中间居住，我要在你们中间行走；我要做他们的天主，他们要作我的子民。所以，你们要从他们中间出来，彼此分离，不要触摸不洁之物。}宗徒啊，你也继续编造这个论调，正当你亲自向如此巨大的不洁漩涡伸出手时；甚至，你更进一步说：\InnerQuote{我们既有这恩许，亲爱的，就让我们从一切肉身和灵魂的玷污中洁净自己，在敬畏天主中成全圣洁。}}请问，那位在我们心中刻下这些（劝诫）的人，难道是在把一个臭名昭著的淫乱者召回到教会里吗？或者，他写这些信的理由，是为了防止在今天你们眼中他好像曾这样做过？这些（他的话语）将同样有义务作为前文（部分）的规范原则，并作为后文（部分）的预判。因为他在书信的结尾说：\BibleQuote{恐怕我来到的时候，天主使我谦卑，我要为许多从前犯罪、而没有为他们所行的不洁、淫乱和放荡悔改的人哀哭}，他当然不是判断说，如果他们进入悔改之路，而那些（他要在教会里找到的人）就要被接回（被他\Emph{进入}教会），而是说他们要被哀哭，并且毫无疑义地被驱逐，好使他们失去悔改（的益处）。而且，此外，他先前已经断言光明与黑暗、义与不义之间没有相通，现在却在此处指出一些有关相通的事，这是不协调的。但所有那些按违反这人本身的自然和意图、违反他教义的规范与法则来理解任何事的人，都是不认识宗徒的；所以他们竟敢假定他，这位一切圣德的导师，甚至以自身为例，一切不洁的谴责者和净化者，在这些方面完全自洽的人，居然会恢复教会特权给一个乱伦的人，而不是给某个更轻微的犯法者。\par

% structure: p0047 source=h2 target=section reason=relative-heading-rank; base=h2

\section{第十六章 宗徒的一贯性}

因此，必须不断地向他们指出宗徒的（品性）；我将凭我在他所有书信中对他的认识，坚持他在\BibleBook{格林多}{后书}中也是如此的。正是他，即便在\BibleBook{格林多}{前书}中，也是所有（宗徒）中第一个将天主的宫殿奉献出来的人：\BibleQuote{难道你们不知道，你们是天主的宫殿，天主圣神住在你们内吗？}——同样，为了祝圣和洁净那宫殿，他撰写了有关宫殿看守者的法律：\BibleQuote{如果有人毁坏了天主的宫殿，天主必要毁坏他，因为天主的宫殿是圣的，这宫殿就是你们。}那么，请问：世上谁曾重新整合过一个被天主\Quote{毁坏}的人（即交付于撒殚以毁灭肉身的）？而在附加了\BibleQuote{谁也不要自欺}之后，即不要有人以为被天主\Quote{毁坏}的人还能重新整合？同样，在所有其他罪恶中——不，甚至是在一切罪恶\Emph{之前}——当他肯定\BibleQuote{奸淫者、行淫者、娈童者、同性交合者，不能承受天主的国}时，他先说了\BibleQuote{不要错误}——即如果你们以为他们能承受。但对于那些被夺去\Quote{天国}的人，天国中的生命自然也不允许。此外，他加上：\BibleQuote{但你们从前确实是这样；但你们已经洗净了，已经圣化了，因主耶稣基督之名，并因我们天主的圣神。}就他而言，将这些罪过放在账目的已付方，即圣洗\Emph{之前}；而在圣洗\Emph{之后}，他判定它们是不可赦免的，如果真是这样（也确实如此）：他们不允许重新\Quote{接受洗净}。也要在接下来的内容中认出\Person{保禄}（作为）纪律及其规范的不可动摇的柱石：\BibleQuote{食物是为肚腹，肚腹是为食物；但天主使这两样都归于无有；然而身体不是为行淫，而是为主。}因为天主说：\BibleQuote{让我们造人，}\BibleQuote{按照我们的肖像和模样。}\BibleQuote{天主造了人；按照天主的肖像和模样造了他。}\BibleQuote{主是为身体；}是的，因为\BibleQuote{圣言成了血肉。}\BibleQuote{此外，天主既使主复活，也必要藉自己的能力使我们复活；}当然，是因我们的身体与祂的结合。因此，\BibleQuote{难道你们不知道你们的身体是基督的肢体吗？}因为\Person{基督}也是天主的宫殿。\BibleQuote{拆毁这座圣殿，我三天内要把它重建起来。}\BibleQuote{我拿基督的肢体作妓女的肢体吗？难道你们不知道，与妓女结合的，便成为一体？因为那两个人要成为一体；但结合于主的，便是与主成为一神。你们要逃避淫乱。}如果藉着宽恕可以撤销，那么我为什么要逃避它，再去重新犯奸淫呢？如果我逃避了，我将一无所获：我将成为\Quote{一个身体}，我因共融而与之结合。\BibleQuote{人所犯的任何罪都在身体以外；但行淫的人是得罪自己的身体。}而且，恐怕你会因为那个说法而跑去行淫，借口你是在得罪你自己的东西，而不是主的，他把你从你自己那里夺走，按照他先前的安排，将你属于\Person{基督}：\BibleQuote{你们不是自己的；}立刻反对说：\BibleQuote{因为你们是被高价买来的}——即主的血：\BibleQuote{要在你们身上光荣和赞美主。}看看，颁布这道命令的人，是否可能赦免了一个羞辱了主、并把主从自己身体的（帝国）中赶出去的人，而且这确实是通过乱伦。如果你愿彻底吸收宗徒的全部知识，以理解他用怎样的审查之斧砍伐、拔除和消灭一切情欲的森林，生怕让任何东西重新萌芽生长；请看，他渴望灵魂斋戒，远离自然的合法果实——我指的是婚姻的苹果：\BibleQuote{关于你们所写的事，男人不接触女人是好的；但是，为了避免淫乱，每个男人应各有自己的妻子；丈夫对妻子、妻子对丈夫，应尽本分。}谁不知道他是违反自己的意愿放宽了这个\Quote{好}的约束，以防止淫乱？但是，如果他曾或正对淫乱予以宽恕，那么他自然破坏了自己补救措施的目的，并且必须立即对节制的婚姻加以限制，如果因之而允许那些婚姻的淫乱不再被畏惧的话。因为被予以宽恕的淫乱不会被畏惧。然而，他承认他准许婚姻的使用是\Quote{出于宽恕，不是出于命令。}因为他\Quote{\Emph{愿意}}所有人与他一样。但当合法的事只是出于宽恕被准许时，谁还指望非法的事呢？他对\Quote{未婚者}和\Quote{寡妇}说：\BibleQuote{照他的榜样，坚持（现状）是好的；}\BibleQuote{但如果他们软弱，就可结婚；因为结婚比燃烧更好。}请问，是比什么火更好\Quote{燃烧}——是贪欲之火，还是惩罚之火？不，如果淫乱是可赦免的，它就不会成为\Emph{贪欲}的对象。但宗徒更倾向于预防\Emph{惩罚}之火。因此，如果是\Emph{惩罚}所指的\Quote{燃烧}，那么等待\Emph{惩罚}的淫乱就是不可赦免的。同时，在禁止休妻时，他用主反对奸淫的诫命作为工具，以提供替代休妻的，要么是寡妇的坚持，要么是和好的和解：因为\BibleQuote{凡休妻的（除非为了奸淫的缘故），就是使她犯奸淫；而娶被丈夫休了的妇人的，也是犯奸淫。}圣神提供了多么强有力的补救措施，以防止他不再愿意赦免的罪行再次发生！\par

如今，若在一切情况下，他说人这样是最好的：\BibleQuote{你与妻子结合，就不要寻求分离}\BibleRef{格前 7:27}（以免你给奸淫提供机会）；\BibleQuote{你从妻子得了释放，就不要寻求妻子}\BibleRef{格前 7:27}，好为你自己保留机会；\BibleQuote{然而，你若娶了妻子，并且处女若出嫁，她并不犯罪；不过，这样的人肉身将受苦难}\BibleRef{格前 7:28}——即使在此，他也是藉着\BibleQuote{宽免他们}\BibleRef{格前 7:28}而给予许可。另一方面，他规定\BibleQuote{时期已缩短}\BibleRef{格前 7:29}，为叫甚至\BibleQuote{有妻子的，要像没有一样}\BibleRef{格前 7:29}。\BibleQuote{因为这世界的局面正在逝去}\BibleRef{格前 7:31}——这世界显然不再需要\BibleQuote{你们要生长和繁殖}\BibleRef{创 1:28}的命令。因此，他愿意我们\BibleQuote{毫无挂虑}\BibleRef{格前 7:32}地度日，因为\BibleQuote{没有妻子的，挂虑主的事，想怎样悦乐天主；娶了妻子的，却挂虑世上的事，想怎样悦乐自己的配偶}\BibleRef{格前 7:32-33}。如此，他宣称\BibleQuote{保守童女的}\BibleRef{格前 7:38}比\BibleQuote{将她出嫁的}\BibleRef{格前 7:38}做得\BibleQuote{更好}\BibleRef{格前 7:38}。同样，他也有区别地判断，那在信仰后失去丈夫的妇女，爱慕地拥抱寡居的机会，是更有福的。因此，他推崇所有这些节制的劝谕为属神的：\BibleQuote{我想，}他说，\BibleQuote{我也具有天主的圣神}\BibleRef{格前 7:40}。\par

这个极其放肆地断言一切不洁的人是谁？显然是奸淫者、淫乱者、乱伦者的\Quote{最忠信}的辩护人，他为此荣誉反对圣神而承担了这项事业，以致他从（祂的）宗徒（的著作）中援引了虚假的见证。保禄没有给予如此宽纵，他努力将\Quote{肉身的需要}从甚至正当的（婚姻结合）借口（列表）中完全抹去。他确实给予\Quote{宽纵}，我承认——不是给予奸淫，而是给予婚姻。他确实\Quote{宽免}，我承认——是婚姻，而非卖淫。他试图避免甚至宽免本性，恐怕他会谄媚罪过。他热心限制那承受祝福的结合，恐怕那承受诅咒的会被原谅。留给他（唯一）的可能性是——从（本性的）渣滓中净化肉身，因为他不能（将肉体）从（肮脏的）污点中洗净。但这正是悖逆无知异端者的惯常做法；是的，到如今甚至对普世属血气者也是如此：用某个模棱两可的段落及时的支持来武装自己，反对整部文献纪律严明的众多句子。\par

% structure: p0051 source=h2 target=section reason=relative-heading-rank; base=h2

\section{第十七章　宗徒在其他书信中的一致性}

请向我挑战，让我面对宗徒的战阵：看看他的书信，它们全都守卫着端庄、贞洁和圣洁，全都将标枪瞄准奢华、放荡和情欲的营垒。简言之，他给得撒洛尼人写了什么？\Quote{因为我们的安慰（并非源于）诱惑，也不是出于不洁：}并且，\Quote{这是天主的旨意，你们的圣化：你们要戒避淫乱；各人当知道怎样以圣化和尊贵持守自己的身体，不要像那些不认识天主的外邦人一样，沉溺于情欲的贪恋。}迦拉达人读到了什么？\Quote{肉身的作为是显而易见的。}这些是什么？在最先列出的中，他放置了\Quote{淫乱、不洁、放荡：}\Quote{（关于这些），我预先告诉你们，正如我已经告诉过你们：行这样事的人，不能继承天主的国。}此外，罗马人——有什么教导比这更深刻地印在他们心中：相信后不可离弃主？\Quote{那么，我们怎么说呢？我们要留在罪中，好使恩宠增加吗？断乎不可。我们这些死于罪的人，怎能还在罪中活着呢？难道你们不知道，我们受洗归于基督耶稣的人，就是受洗归于祂的死吗？所以我们藉着洗礼归于死，与祂同葬，好使基督怎样从死者中复活，我们也怎样在新生活中行走。因为我们如果与祂同死似的形状结合，那么，也定要与祂同复活；知道这一点：我们的旧人已与祂同钉在十字架上。但如果我们与基督同死，我们相信也要与祂同活；知道基督既从死者中复活，就不再死，死不再统治祂。因为祂死，是死于罪，\Emph{一次而竟全功}；祂活，是活于天主。同样，你们也要看自己是死于罪，但在基督耶稣内活于天主。}因此，基督既一次死去，那么在基督之后死去的人，不能再活于罪，尤其是如此可憎的罪。否则，如果淫乱和奸淫可能再次被容许，那么基督也就有可能再次死去。此外，宗徒迫切地禁止\Quote{罪在我们必死的身体中为王，}他知道这身体的\Quote{肉身的软弱。}\Quote{因为你们从前怎样将肢体献给不洁和不法作奴仆，现今也要照样将肢体献给正义，成就圣德。}因为他虽已肯定\Quote{善不住在他的肉身内，}但这是指\Quote{文字的律法，}他曾在这律法之下；但按照他使我们依附的\Quote{圣神的律法，}他解救我们脱离\Quote{肉身的软弱。}他说：\Quote{因为生命之圣神的律法，已使你从罪与死的律法中获得了自由。}尽管他可能看起来部分是以犹太主义的立场辩论，但他是指向我们，把纪律规则的完整性与圆满性指示给我们——我们这些曾在律法下劳苦的人，正是为了我们，\Quote{天主派遣了自己的儿子，带着罪恶肉身的形状，为了罪恶，并在肉身中定了罪案，}他说，\Quote{为使律法的正义，成就在我们这些不随从肉身，而随从圣神的人身上。因为随从肉身的人，思念肉身的事；随从圣神的人，思念圣神的事。}此外，他肯定\Quote{肉身的思念}是\Quote{死亡，}因此也是\Quote{敌对，}且是敌对\Emph{对天主}；并且\Quote{那些在肉身内的人，}即是在肉身的\Emph{思念}里，\Quote{不能悦乐天主：}他又说：\Quote{如果你们随从肉身生活，必要死亡。}但我们明白\Quote{肉身的思念}和\Quote{肉身的生活}指的是什么，除非是那些\Quote{羞于启齿}的事？因为肉身的其他作为，就连宗徒也曾提名道姓。同样，在写给厄弗所人的信中，他在回忆过去时，告诫他们未来：\Quote{我们从前也在其中生活，放纵肉体的私欲和享乐。}最后，他指责那些背弃自己的人——即基督徒——因为他们\Quote{将自己交付给各种不洁的行为，}\Quote{但你们，}他说，\Quote{却不是这样学了基督。}他又如此说：\Quote{从前偷窃的，不要再偷。}但同样，从前惯于奸淫的，不要再奸淫；从前惯于淫乱的，不要再淫乱：因为他若习惯宽恕这样的人，或愿意宽恕他们，他也会加上这些告诫——他不愿连一句话也玷污人，故说：\Quote{不可出口败坏的话。}又说：\Quote{至于淫乱和一切不洁，在你们中连提也不可提，这样才合圣人的身份——}这远非可以原谅——\Quote{你们知道，凡是淫乱或不洁的人，在天主的国里没有产业。没有人用虚言欺骗你们，正是为了这些事，天主的忿怒才降在悖逆之子身上。}谁\Quote{用虚言欺骗}？岂不是那在公开宣讲中声称奸淫可以宽恕的人？他没有看见宗徒已经挖掉了它的基础，当他限制酗酒和狂欢时，也在这里说：\Quote{不要醉酒，醉酒使人放荡。}他也给哥罗森人指出，他们在地上应当\Quote{致死}哪些\Quote{肢体：}\Quote{淫乱、不洁、邪情、恶欲}和\Quote{污秽的言语。}到此为止，面对如此众多而明确的经文，请你放弃你所执着的唯一段落吧。少数被多数掩盖，疑惑被确定淹没，隐晦被清晰淹没。即使宗徒确实曾宽恕了那位格林多人的淫乱，那也是他为了适应时代需要，一次性违反自己常例的另一个例子。他只给弟茂德行了割损，却废除了割损。\par

% structure: p0053 source=h2 target=section reason=relative-heading-rank; base=h2

\section{第十八章 答属魂派的异议。}

\Quote{但这些（经文）}，（我们的对手）说，\Quote{将属于禁止一切不贞洁，并强制一切端庄，但无损于赦免之处；这（赦免）在罪被定罪时并非立即完全被否认，因为赦免的时间与它所排除的定罪同时发生。}\par

属魂派方面的这一机智部分是（自然地）连贯的；因此我们为此处保留了那些警诫，甚至在古代，这些警诫也公开地采取，以拒绝将这类情况纳入教会的共融。\par

因为甚至在\BibleBook{}{箴言}中，我们称之为\OriginalTerm{箴言}{Parœmiæ}，撒罗满特别（论及）奸夫（是）在任何地方都不容于补赎的。\BibleQuote{但奸夫}，他说，\BibleQuote{因缺乏理智，为自己的灵魂招致丧亡；忍受痛苦和耻辱。况且，他的耻辱永不被抹去。因为充满嫉妒的忿怒，在审判之日必不饶恕那人。}你若以为这是论及外邦人，无论如何，论及信徒，你已通过\BibleBook{}{依撒意亚}听到（这话）：\BibleQuote{你们从他们中间出来，分别出来，不要触摸不洁之物。}你在\BibleBook{}{圣咏}的开头就听到：\BibleQuote{凡不随从恶人的计谋，不站罪人的道路，不坐亵慢人的座位，这人是有福的；}随后他的声音（也出现）：\BibleQuote{我没有与虚妄的集会同坐，也不与行不义的人同进。}——这（涉及）那些行恶之人的\BibleQuote{\Emph{教会}}——\BibleQuote{不与不虔敬的人同坐；}并且，\BibleQuote{我要洗手表明无辜，好环绕祢的祭坛，上主}——作为\Quote{独自一人的军队}——因为确实\BibleQuote{与圣洁的人，祢显为圣洁；与纯全的人，祢显为纯全；与蒙拣选的人，祢显为蒙拣选；与乖僻的人，祢显为乖僻。}并且在别处：\BibleQuote{但上主对罪人说：你为何传述我的正义，你的口又提及我的盟约？你若看见盗贼，就与他同伙；又与奸夫同分。}因此，宗徒从这里得到教训，也说：\BibleQuote{我在书信中写信给你们，不可与淫乱的人相交：当然不是指这世上的淫乱的人}——等等——\BibleQuote{否则你们必须离开这世界。但如今我写信给你们，若有人被称为弟兄，是淫乱的，或拜偶像的}（因为什么如此亲密？），\BibleQuote{或欺诈的}（因为什么如此相近？），等等，\BibleQuote{这样的人连一起吃饭也不可}，更不用说领圣体了：因为，诚然，\BibleQuote{一点面酵能使全团发起来。}又对弟茂德说：\BibleQuote{不可仓促给人覆手，也不可参与别人的罪过。}又对厄弗所人说：\BibleQuote{不要与他们同伙；因为你们从前是黑暗。}并且更恳切地说：\BibleQuote{不要参与黑暗无益的作为；倒要责备它们。因为他们在暗中所行的，就是提起来也是可耻的。}有什么比不贞洁更可耻的呢？再者，即便从\BibleQuote{弟兄}中，他警告得撒洛尼人要远离那些\BibleQuote{游手好闲}的人，何况从淫乱的人呢！因为这些都是基督的定论，\BibleQuote{祂爱教会，}并\BibleQuote{为自己舍弃了她，为圣化她（用水洗净她，藉着圣言），好使祂把教会呈献给祂自己，光荣的，没有斑点或皱纹}——当然是在洗净\Emph{之后}——\BibleQuote{而为圣洁无瑕；}此后，即是\BibleQuote{无皱纹}如童贞女，\BibleQuote{无斑点}（无淫乱之玷）如新娘，\BibleQuote{无耻辱}（无卑劣），因为已彻底\BibleQuote{洗净}。\par

假若在此，你仍然想要回答说：圣体共融确实被拒绝给予罪人，尤其是那些曾被\Quote{肉身所玷污}的人，但只是暂时的；即，藉补赎的恳求而得以恢复——这是依照天主的仁慈，祂宁愿罪人悔改胜过其死亡？——因为必须普遍地攻击你们意见的这一根本基础。为此我们说，如果天主的仁慈能够保证其彰显甚至及于领圣洗后失足的人，那么宗徒本应这样说：\Quote{不要与黑暗的行为共融，\Emph{除非他们悔改}}；以及\Quote{连与这样的人一起进食也不可，\Emph{除非他们先用脚滚地擦干弟兄们的鞋子}}；以及\Quote{谁若摧毁天主的殿，天主必摧毁他，\Emph{除非他在教会中从自己头上抖落所有炉灶的灰烬}}。因为他的职责是，对于那些他所谴责的事，同样地确定他谴责（并且有条件地谴责）的程度——他是否以暂时和有条件而非永久的严厉谴责了它们。然而，既然在所有书信中，他都禁止此等（在相信后仍如此犯罪）的人被接纳（入信徒团体）；而若被接纳，则将之从共融中逐出，没有任何条件或时间的希望；他更赞同我们的意见，指明主所喜爱的悔改，是在相信之前、在圣洗之前，被认为优于罪人之死的悔改——（罪人，我说）藉基督的恩宠一次而永远地被洗涤，基督为我们的罪一次而永远地死了。因为宗徒在他自己身上也立下了这一（规则）。当他肯定基督为此目的而来，即祂可能拯救罪人——他自己曾是他们中的\Quote{首恶}——时，他又加了什么？\Quote{我蒙受了怜悯，因为我在不信中无知地做了这事}。因此，天主的仁慈——宁愿罪人悔改胜过其死亡——是针对那些仍然无知、仍然不信的人，基督为拯救他们而来；而非针对那些已认识天主、已学会信仰的圣事的人。但如果天主的仁慈适用于仍然无知和不信的人，那么当然，悔改将仁慈引向自身；这并不损害相信之后的悔改类型，这种悔改对于较轻的罪，可以从主教那里获得宽恕，对于较重且不可饶恕的罪，则只能从天主那里获得。\par

% structure: p0058 source=h2 target=section reason=relative-heading-rank; base=h2

\section{第十九章　驳斥从\BibleBook{}{默示录}和\BibleBook{若望}{一书}提出的反对意见}

但我们（要处理）保禄到何种程度呢？因为甚至若望似乎也暗中给对立面某种支持。仿佛在\BibleBook{}{默示录}中，他显然为奸淫指定了悔改的辅助，当圣神向提雅提辣的天使传达信息说，祂\Quote{反对他，因为他容忍那妇人依则贝耳——她自称是先知的，教导并诱惑我的仆人行奸淫和吃祭偶像之物。我曾宽限她时日，让她悔改，但她不愿因奸淫而悔改。看，我要把她投入病床，并使与她行奸淫的人同她一起遭受大患难，除非他们悔改她的作为。}我满足于宗徒之间在信仰和纪律的规则上有共同约定。因为\Quote{无论是我，}保禄说，\Quote{还是他们，我们都如此宣讲。}因此，关乎整个圣事的利益，相信若望允许了任何被保禄明确拒绝的事，这是不紧要的。谁若留意圣神的这种和谐，必被祂引导进入祂的深意。因为（提雅提辣教会的天使）正秘密地将一个异端妇人引入教会，并恰当地催促她悔改，这妇人曾承担教导她从尼苛劳党人所学的内容。因为谁怀疑一个异端者，被（虚假的圣洗）仪式所欺骗，在发现他的不幸后并藉悔改赎罪，既获得赦免，又被恢复进入教会的怀抱？因此，甚至在我们中间，正如与一个外邦人同等，甚至超过外邦人，一个异端者同样被真理的圣洗从双重身份中被洁净，并被接纳（入教会）。否则，如果你确信那妇人是在活信德之后后来又死了并成为异端者，以便你主张悔改的结果是赦免——不是为异端者，而是为相信的罪人：我承认，让她悔改吧；但目的是停止奸淫，而不是为了恢复（教会共融）。因为这将是一种我们同样承认（比你）更应归于其的悔改；但我们为赦免把它保留给天主。\par

简言之，这部\BibleBook{若望}{默示录}在后来的段落中，将\BibleQuote{可憎者和行淫者}以及\BibleQuote{胆怯的、不信的、可憎的、杀人的、行邪术的、拜偶像的}——那些在明认信仰时犯下任何此类罪行的人——判归\BibleQuote{火湖}，没有任何\Emph{有条件性的}判罪。因为这不像是针对\Emph{外邦人}的（判词），因为它（刚刚）论及\Emph{信徒}时说：\BibleQuote{得胜的，必承受这些为产业；我要作他们的天主，他们要作我的子女；}接着又说：\BibleQuote{但胆怯的、不信的、可憎的、行淫的、杀人的、行邪术的、拜偶像的，他们的分就在烧着硫磺的火湖里；这是第二次的死。}同样，又说：\BibleQuote{遵守诫命的人是有福的，可得权柄到生命树那里，也能从门进城。狗、行邪术的、行淫的、杀人的、偶像崇拜者，都出去！}——当然，是指那些\Emph{不}遵守诫命的人；因为\Emph{被赶出去}是那些\Emph{曾在里面的人}的分。此外，\BibleQuote{审判教外的人与我何干？}这话已经先于（上述判决）说了。\par

他们立刻从\BibleBook{若望}{一书}中摘取（一个证据）。经上说：\BibleQuote{他圣子的血洗净我们一切的罪。}那么，如果祂常常且从一切罪中洗净我们，我们就会常常以各种形式犯罪；否则，如果不是\Emph{常常}，那么就不是在信仰之后再次（犯罪）；如果不是从罪中（洗净），那么就不是从行淫中（洗净）。但是（若望）从何处入手呢？他宣认\BibleQuote{天主}是\BibleQuote{光}，且\BibleQuote{在他内毫无黑暗}，又\BibleQuote{如果我们说我们与他相通，却在黑暗中行走，就是说谎}。他说：\BibleQuote{如果我们却在光中行走，我们就与他相通，而他圣子耶稣基督的血会洗净我们一切的罪。}那么，在光中行走，我们会犯罪吗？在光中犯罪，我们会被彻底洗净吗？绝不。因为犯罪的人不是在光中，而是在黑暗中。由此他也指出了我们将如何彻底洗净罪的方式——就是\BibleQuote{在光中行走}，在光中罪不能犯。因此，他说我们\Quote{被彻底洗净}的意思，不是因着我们犯罪，而是因着我们\Emph{不}犯罪。因为\Quote{在光中行走}，不与黑暗相通，我们就会像那些\Quote{被彻底洗净}的人一样行事；罪虽未完全根除，但不再明知故犯。这就是主宝血的能力：它既已洗净罪，便从此将人安置在\Quote{光中}，并使他们从此纯洁，只要他们持续在光中行走。你说：\Quote{但他又补充说：\BibleQuote{如果我们说我们没有罪，便是自欺，真理不在我们内。如果我们承认我们的罪，他是忠信正义的，必赦免我们的罪，并洗净我们一切的不义。}}他说的是\Quote{不洁}吗？不是。否则，如果是这样，那么他也\Quote{彻底洗净}我们\Quote{拜偶像}的罪了。但意义上是有区别的。再看：他说：\BibleQuote{如果我们说我们没有犯过罪，便是以他为说谎者，他的话不在我们内。}更全面地说：\BibleQuote{孩子们，我写这些给你们，是要你们不犯罪；但若有人犯了罪，我们在父那里有一位护慰者，就是正义的耶稣基督；他为我们作了赎罪祭。}你说：\Quote{根据这些话，我们必须承认我们既犯罪，又得赦免。}那么，当我继续读这书信，发现不同的话时，你的理论又当如何呢？因为他断言\Emph{我们根本不犯罪}；并且为此详加论述，以免作出这样的让步；阐明罪已被基督一次性消除，此后不再获得赦免；在这一点上，上下文要求我们将这论述应用于对\Emph{洁德}的劝勉。他说：\BibleQuote{凡怀着这希望的，必使自己贞洁，如同他一样贞洁。凡犯罪的，也就行不法；罪就是不法。你们知道，他显现出来，是为除掉罪}——从此以后当然就不再犯罪了，如果他（紧接着）说：\BibleQuote{凡住在他内的，就不犯罪；凡犯罪的，是没有看见他，也不认识他。孩子们，不要让人迷惑你们。凡行正义的，就是正义的，如同他一样正义。凡犯罪的，就是属于魔鬼，因为魔鬼从起初就犯罪。天主子所以显现出来，是为消灭魔鬼的作为。}因为他已经\Quote{消灭}了它们，通过圣洗使人获得自由，将\Quote{死亡的债务}\Quote{当作恩赐}赐给了他：因此，\BibleQuote{凡由天主生的，就不犯罪，因为天主的种子存留在他内；他不能犯罪，因为他是由天主生的。这样，天主的子女和魔鬼的子女就显露出来了。}如何显露呢？除非是这样：前者因不犯罪，自从由天主而生时起；后者因犯罪，因为他们属于魔鬼，就好像从未由天主而生一样？但如果他说：\BibleQuote{不正义的人不是属于天主的}，那么那已不再属于天主、不\Emph{端庄}的人，又怎能再次成为天主的（儿子）呢？\par

\Quote{因此这几乎等于说\Person{若望}忘记了自己：在他的书信前半部分断言我们并非没有罪，现在却又规定我们根本不可犯罪；一方面以赦免的希望稍加奉承我们，另一方面却严厉断言，凡犯罪者都不是天主的子女。}但是请除去（这种想法）：因为我们自己也没有忘记罪之间的区别，这正是我们离题开始之处。（这是一个正确的区别），因为\Person{若望}在此认可了它；因为有些罪是日常所犯，我们都可能犯：因为谁能免于以下意外：或因不义而发怒，且怒气过日落；或甚至动手施暴；或言语不慎中伤人；或轻易发誓；或背弃承诺；或出于羞怯或\InnerQuote{必要？}而说谎？在事务中，在职务中，在买卖中，在饮食中，在观看中，在听闻中，我们受多大的诱惑摆布！因此，如果这些罪没有宽恕，任何人都无法得到救恩。对于这些，藉着圣父成功的转求者基督，将会有宽恕。但也有与此相反者，如更严重且具有毁灭性的罪，不能获得宽恕——凶杀、偶像崇拜、欺诈、背教、亵渎；当然也包括奸淫和淫乱；以及任何其他\InnerQuote{侵犯天主殿堂}的行为。对于这些，基督将不再是成功的辩护者：凡由天主而生者必不会犯这些罪，他若犯了这些罪，就不再是天主的子女。\par

因此，\Person{若望}关于多样性的规则得以确立；他如此安排罪的区别，时而承认，时而否认天主的子女犯罪。因为（在作这些断言时）他正望向其书信的最终条款，并为那（最终条款）奠立初步基础；意在结尾更明确地说：\BibleQuote{若有人看见自己的弟兄犯了不至于死的罪，就应当祈求，主必将生命赐给那犯罪不至于死的人。有至于死的罪，我不说要为那罪祈求。}他（像我一样）也记得，天主曾禁止\Person{耶肋米亚}为代表犯大罪的百姓求情。\BibleQuote{一切不义都是罪，但也有不至于死的罪。但我们知道，凡由天主生的，就不犯罪}——即那至于死的罪。因此，你只剩下一个选择：要么否认奸淫和淫乱是大罪，要么承认它们不可赦免，因为连为它们成功的代祷也不被允许。\par

% structure: p0064 source=h2 target=section reason=relative-heading-rank; base=h2

\section{第二十章：\Person{戴尔都良}从宗徒教导转向宗徒同伴和法律。}

因此，宗徒的纪律（如此称呼）确实教导并坚定地指导，作为要点，关于天主殿堂的圣洁监督者，要彻底根除对贞洁的一切亵渎侵犯，丝毫不提恢复。但我还希望额外补充一位特定宗徒同伴的见证——这见证最适合以最近的权利来证实其导师的纪律。因为有一封以\Person{巴尔纳伯}之名流传的\WorkTitle{希伯来人书}——他是天主充分认可的人，正如\Person{保禄}在持续持守节欲中将他安置在自己旁边：\BibleQuote{除了我和\Person{巴尔纳伯}，难道我们没有工作的权柄吗？}当然，\WorkTitle{巴尔纳伯书信}在教会中比那奸淫者的伪经\WorkTitle{牧人篇}更被普遍接受。因此，他警告门徒要离开一切初步的道理，努力追求成全，不再奠立从死行中悔改的根基，他说：\BibleQuote{因为那曾经被光照，尝过天主的恩赐，又分享了圣神，并尝过天主圣言的甘甜，但当他们的年岁已暮而行将堕落时，不可能重新唤起他们悔改，为自己再次钉死天主子，并羞辱祂。}\BibleQuote{因为土地吸收了常降的雨水，并生出合用的蔬菜，得享天主的祝福；但若长出荆棘，它便被摈弃，近于诅咒，它的结局就是焚烧。}他\Emph{从}宗徒学习这些，并\Emph{与}宗徒一同教导，从不曾知道任何宗徒应许给奸淫者和淫乱者的\Quote{第二次悔改}。\par

因为他善于解释法律，即使在真理本身（的计划）中仍保持其预像。简言之，正是参照这种纪律，才针对麻风病人采取了谨慎措施：\Quote{但若斑点已在皮上开花，盖满了全身，从头到脚，凡肉眼所见之处，那么司铎一见，就应完全洁净他；因为他已完全变白，是洁净的。但在那一日，若他身上出现了红润的肉，他就不洁净。}律法的意思，是那从肉体的原始状态完全转变为信德之洁白的人——这信德在世人的眼中被视为缺陷和瑕疵——并且完全成为新造的人，应被理解为是\Quote{洁净的}；因为不再是\Quote{有斑点的}，不再掺杂原始与新生的东西。然而，如果在（不洁判决的）撤销之后，旧本性的任何部分连同其倾向重新活跃起来，那么那曾在他肉身中被认为对罪完全死了的，必须再次被视为不洁，并且不能再由司铎行赎罪礼。因此，奸淫从原始根苗重新发芽，完全玷污了那曾被排除的新颜色的合一，是一种无法洁净的瑕疵。再者，就房屋而言：如果有人在隔墙上发现任何斑点和凹坑，报告给司铎，在他进去察看那房屋之前，他吩咐把所有（室内物品）都搬走；这样，房中的物品就不致不洁。然后，司铎进去后，若发现青绿色或淡红色的凹坑，且其外观在隔墙体内显得更深，他要出去到门口，将那房屋封锁七天。然后，到第七天回来时，若他发现那污点已在隔墙上扩散，他就要下令将那沾染了癞病的石头挖出来，丢在城外不洁的地方；然后取别的磨光而完好的石头，放在原先的石头处，并用别的灰泥粉刷房屋。因为，当人来到圣父的大司祭——基督面前时，必须先在一周内除去一切障碍，使那存留的房屋，即肉身和灵魂，成为洁净；及至天主圣言进入其中，发现\Quote{红与绿的污点}时，就必须立即将那致命而嗜血的激情\Quote{挖出来}并\Quote{丢到}门外——因为\WorkTitle{默示录}也将\Quote{死亡}置于\Quote{绿马}上，而将\Quote{战士}置于\Quote{红马}上——然后，要用磨光而适于衔接的坚固石头铺在它们的位置上，就是那些由天主使之成为亚巴郎的子孙的石头；这样，人才能适合天主。但是，如果在那房屋复原和改造之后，司铎又在那房屋里发现任何原始紊乱和瑕疵，他就宣布它不洁，并命令将木材、石头和所有结构都拆毁，丢到不洁的地方。这个人就是——肉身和灵魂——在改造之后、领受圣洗和司铎进入之后，又重新生发肉体的疮痂和污点，并\Quote{被丢到城外不洁的地方}——即\Quote{被交付给撒殚，以毁灭肉身}——并且在他覆亡之后，不再在教会中重建。同样，关于与一个已订婚却尚未赎出、尚未释放的女奴同床的事：律法说，要为她\Quote{设法规矩}，她不应被处死，因为她尚未被释放给那为谁保存她的人。因为那尚未被释放给基督的肉——它曾为祂保存——曾经可以污染而免受惩罚；同样，如今在解放之后，它不再得到宽恕。\par

% structure: p0067 source=h2 target=section reason=relative-heading-rank; base=h2

\section{第二十一章 论纪律与权力的区别，以及钥匙的权力}

如果宗徒们更明白这些（法律的比喻意义），他们当然更加谨慎（对待它们，甚至超过宗徒时代的人）。但如今我甚至要降到这争辩的地步，在宗徒的\Emph{教导}与他们的\Emph{权能}之间作出区分。纪律治理人，权能盖印于人；此外，权能是圣神，但圣神是天主。况且，（圣神）曾教导什么？不可与黑暗的行为相通。看祂吩咐什么。再者，谁曾能赦免罪过？这是祂独有的特权：因为\Quote{除了天主以外，谁能赦罪呢？}而且，当然，（除了祂外，谁能赦免）\Emph{大}罪，就是那些得罪祂自己和祂圣殿的罪？因为，就你而言，那些属于对你个人冒犯的罪，你奉（主）命，在伯多禄身上，要赦免直到七十个七次。所以，如果连蒙福的宗徒们也曾准许过任何此类宽免（得罪），其宽恕来自天主而非人，那么他们这样做是可行的，不是凭纪律，而是凭权能。因为他们既复活死者——这只有天主（能行），又使衰弱者恢复完整——这只有基督（能行）；不但如此，他们也降灾祸，而基督所不愿行的。因为祂来受苦，不宜严厉。阿纳尼雅和厄里玛都被打击——阿纳尼雅死于死亡，厄里玛瞎了眼——为借此证明基督也曾有能力行\Emph{同样}的（奇迹）。同样，（古时的）先知们也曾赐予悔改者杀人的\Emph{赦免}，连同奸淫，因为他们同时显示了\Emph{严厉}的明显证据。因此，现在也请向我展示，宗徒式的主教啊，先知的证据，使我认识你的神圣德能，并为你自己辩护那赦免这类罪的\Emph{权能}！然而，如果你只被授予了\Emph{纪律}的职分，以及（职责）不是王者式的而是服务性的管理权，那么你是谁或你何其大，竟能赐予宽恕？你既不具先知也不具宗徒的特征，就缺乏那专属于宽恕的德能！\par

\Quote{但是}，你说，\Quote{\Emph{教会}有权赦免罪过。}这一点我承认并更（比你）断定，我拥有圣神自身在新先知们身上，说：\Quote{教会有权赦免罪过；但我不愿如此行，免得他们更犯别的罪。}\Quote{假如是假先知的神说了这话呢？}不然，那破坏者一方面更会以宽仁自夸，另一方面影响其他人去犯罪。或者，如果（假先知的神）是依据\Quote{真理之神}急于表达此（观点），那么\Quote{真理之神}确实有\Emph{权能}宽免地赦免犯奸淫者，但\Emph{不愿}如此行，如果这会给多数人带来祸害。\par

我现在探问你的意见，（看看）你从何处篡夺这\Quote{教会}的权利。\par

如果因为主曾对伯多禄说：\Quote{在这磐石上，我要建立我的教会，}\Quote{我把天国的钥匙交给了你；}或\Quote{凡你在地上所捆绑或释放的，在天上也要被捆绑或释放，}你就因此以为捆绑和释放的权能已传给了你，即传给了每个与伯多禄有亲缘的教会，那么你是怎样的人呢？你颠覆并全然改变了主明白的意向，祂（的意向）是将此（恩赐）亲自赐予伯多禄的。祂说：\Quote{\Emph{在你身上}，我要建立我的教会；}以及\Quote{我要把钥匙交给\Emph{你}，}而不是\Emph{教会}；并且\Quote{\Emph{你所释放或捆绑的}一切，}而不是\Emph{他们}所释放或捆绑的。因为结果也是这样教导的。在（伯多禄）身上教会得以建立，即\Emph{通过}（伯多禄）本人；（伯多禄）本人试用这钥匙；你看\Emph{哪}把钥匙：\Quote{以色列人哪，请听我的话：纳匝肋人耶稣，是一位由天主为你们所选定的人，}等等。因此，（伯多禄）本人是第一个在基督的洗礼中打开天国入口的，在那（天国）里面，先前被\Quote{捆绑}的罪得以\Quote{释放}；而那些未被\Quote{释放}的仍被\Quote{捆绑}，符合真正的救恩；他\Quote{捆绑}了阿纳尼雅以死亡的束缚，并\Quote{释放}了脚弱之人脱离健康缺陷。此外，在关于是否遵守法律的争论中，伯多禄是众人中第一个被圣神充满的，他先讲了论及外邦人的召叫，然后说：\Quote{现在你们为什么试探天主，要把我们和我们祖先都不能负的轭放在弟兄们的颈上呢？但是我们相信，我们得救是因主耶稣的恩宠，和他们一样。}这句话既\Quote{释放}了那些被废弃的法律部分，又\Quote{捆绑}了那些被保留的部分。因此，托付给伯多禄的释放和捆绑的权能，与信徒的\Emph{大}罪无关；如果主曾给他一条诫命，要他必须赦免得罪\Emph{他}的弟兄直到\Quote{七十个七次}，那么当然，主会吩咐他此后\Quote{捆绑}——即\Quote{保留}——\Emph{什么}，除非或许有人得罪了\Emph{主}，而不是\Emph{弟兄}。因为得罪\Emph{人}的罪的宽恕，是对得罪\Emph{天主}的罪赦免的一种偏见。\par

现今，这与教会，并与\Emph{你们的}（教会），即\OriginalTerm{属魂的}{Psychic}教会，有何关系？因为按照\Person{伯多禄}的位份，这权柄将相应属于属神的人，无论是宗徒还是先知。因为教会本身，正确而首要地，就是圣神自己，在祂内有一神性之圣三——圣父、圣子、圣神。圣神结合了主所使之成为\Quote{三个}的教会。因此，从那时起，任何结合在这信仰中的数目（的人）都被视为\Quote{一个教会}，出自（教会的）创始者和祝圣者。因此，\Quote{教会}确实会赦免罪过，但那是圣神的教会，借着属神的人，而不是由许多主教组成的教会。因为权柄和审判是属于主的，而非仆人的；属于天主自己，而非司铎的。\par

% structure: p0073 source=h2 target=section reason=relative-heading-rank; base=h2

\section{论殉道者及其为罪大恶极者代祷}

但是，你们甚至将这\Quote{权柄}慷慨地施与殉道者！一旦有人按预先的约定戴上锁链——况且在现今流行的名义拘押下，这些是柔软的锁链——奸淫者就围住他，犯奸淫者接近他；立刻祈祷回响在他周围；立刻所有被玷污者的眼泪（从眼睛）汇成池沼围绕他；没有谁比那些失去（教会共融）的人更热心于购买进入监狱的机会！男人和女人在黑暗中遭侵犯，这黑暗是习惯性纵欲已显然使他们熟悉的；他们从那些冒着自身风险的人手中寻求平安！另一些人投身矿场，然后以领圣体者的身份从那里返回，而此时，对于在\Quote{殉道}之后所犯的罪，需要另一个\Quote{殉道}。\Quote{好吧，世上凡有血肉者谁无罪过？}什么样的\Quote{殉道者}（仍然）是世界的居民，在恳求？手里拿着钱？受制于医生和高利贷者？现在，假设（你们的\Quote{殉道者}）在刀剑之下，头颅已经稳稳地摆好；假设他在十字架上，身体已经伸展；假设他在火刑柱上，狮子已经放出；假设他在轮轴上，木柴已经堆好；在确凿无疑的，我说，确据和拥有殉道之时：谁允许\Emph{人}宽恕那些该保留给\Emph{天主}的（过犯），那些过犯已被天主定罪而不予免除，就连宗徒（就我所知）——甚至殉道者自身——也未曾认为可宽恕？简言之，\Person{保禄}早已\BibleQuote{在厄弗所与野兽搏斗}，当时他判决那乱伦者受\BibleQuote{毁灭}。让殉道者满足于洗净自己的罪过：把以高价获得的也慷慨施与他人，这是忘恩或骄傲的表现。谁曾以自己的死赎买他人的死？唯天主圣子而已。因为在祂的苦难中，祂甚至释放了那个强盗。祂来到世上正是为此：祂自身纯洁无罪，在各方面均圣洁，以便为罪人接受死亡。同样，你们效法祂宽恕罪过，如果你们自己未曾犯罪，显然可以为我受苦。然而，如果你们是罪人，你们那微小火炬的油如何能满足你们和我呢？\par

我现在甚至有一个试验，用以证明基督（在你们内）。若基督在殉道者内，为的是让殉道者可赦免奸淫者和行淫者，那么就让祂公开说出内心的秘密，以便如此赦免罪过；那祂就是基督。因为主耶稣基督正是这样显示祂的权能：\BibleQuote{你们为什么心里思念恶事？为哪一样更容易：对瘫子说\InnerQuote{你的罪赦了}，或是说\InnerQuote{起来行走}？但为叫你们知道，人子在世上拥有赦罪的权柄——我对你说：瘫子，起来，行走。}若主如此看重祂权能的证据，甚至揭示心思，并借此以祂的命令赐予健康，以免祂不被相信有赦罪的权能；那么，没有同样的证据，我就不能相信这同样的权能（存在于）任何人身上——无论他是谁。然而，当你急切地向殉道者恳求为奸淫者和行淫者求赦免时，你自己就承认，这类罪行除非犯罪者本人的殉道，否则不能被洗去，而你却妄图（它们能被）别人的殉道洗去。若果真如此，那么殉道将是另一种圣洗。因为祂说：\BibleQuote{我也有另一种圣洗。}正因如此，从主肋旁的伤口中流出了水和血，即两种圣洗的材料。那么，我也应当藉\Emph{第一}种圣洗拥有释放他人的（权利），如同我藉\Emph{第二}种所能做的；我们必须迫使（我们的对手）得出这个结论：无论是什么权威，什么理由，恢复奸淫者和行淫者与教会的平安，那同样的东西也必然要帮助杀人者和拜偶像者悔改——至少是帮助背教者，当然也包括那在宣认的战斗中，经过与酷刑的艰苦搏斗而被残暴击败的人。此外，若天主和祂的慈悲——祂更愿罪人悔改而非死亡——竟使那些在情欲冲动中跌倒的人比那些在短兵相接中跌倒的人更容易回到教会（的怀抱），这是与天主不相称的。义愤驱使我们说话。你们召回的是被玷污的身体，而不是被血染的身体！哪一种悔改更可怜——是那使被瘙痒的肉身俯伏的悔改，还是那使被撕裂的肉身俯伏的悔改？在所有情况下，哪一种赦免更公正地应被赐予——是自愿的罪人恳求的，还是非自愿的罪人恳求的？没有人\Emph{在}其意愿之下被迫背教；没有人\Emph{违背}其意愿行淫。情欲除自身外不受到任何暴力：它不知道任何强迫。相反地，背教是由何等残忍的巧计和刑罚的种类所强加！谁才是更真实地背教了——那在痛苦中失去基督的人，还是那在享乐中失去基督的人？那在失去祂时忧伤的人，还是那在失去祂时嬉笑的人？然而，那些刻在基督徒斗士身上的伤疤——当然是令基督羡慕的伤疤，因为它们渴求胜利，也因此是荣耀的，因为虽然未能取胜却仍奋战；这些伤疤连魔鬼自己也为之叹息；这些伤疤带有它们自身的不幸，却是贞洁的不幸，带着忧伤的悔改，向主祈求赦免而不羞愧——这些伤疤将被重新归还给这样的人，因为他们的背教是可赎的！惟独在他们身上，\BibleQuote{肉身软弱}。不，没有肉身如此坚强，以致能压碎圣神！\par

\SourceRecord{Translated by S. Thelwall.；Ante-Nicene Fathers；Vol. 4.；Edited by Alexander Roberts, James Donaldson, and A. Cleveland Coxe.；Buffalo, NY: Christian Literature Publishing Co.,；1885.；<http://www.newadvent.org/fathers/0407.htm>.}
